\documentclass[11pt]{article}
\usepackage[T1]{fontenc}
\usepackage{amsmath,amsthm,amssymb,fullpage,longtable,array,microtype,xurl,hyperref}

\emergencystretch=3em
\Urlmuskip=0mu plus 2mu
\let\origunderscore\_
\renewcommand{\_}{\origunderscore\allowbreak}

\newtheorem{theorem}{Theorem}[section]
\newtheorem{lemma}[theorem]{Lemma}
\newtheorem{proposition}[theorem]{Proposition}
\newtheorem{corollary}[theorem]{Corollary}
\theoremstyle{definition}
\newtheorem{definition}[theorem]{Definition}
\theoremstyle{remark}
\newtheorem{remark}[theorem]{Remark}

\title{Erd\H{o}s Problem \#647:\\
A Lean-Verified Structural Reduction\\
and a Conditional Dickson Dichotomy}
\author{Scott D. Hughes}
\date{May 24, 2026}

\begin{document}
\maketitle

\begin{abstract}
Erd\H{o}s Problem~\#647 asks whether any $n>24$ satisfies
$\max_{m<n}(m+\tau(m))\le n+2$. Building on earlier reductions,
we show that any solution has the form $n=2520N$ with $N$ in one
of $96$ residue classes modulo $M=46189=11\cdot 13\cdot 17\cdot 19$.
Under exact-v2 closure semantics (which retires the prior peeled-cofactor
bridge after it was shown unsound), we Lean~4-verify $55$ of these
classes via explicit closure certificates ($42$ single-overlap branch
contradictions and $13$ direct full-value arguments). The remaining
$41$ unresolved residues decompose as the historical hard-$5$
structural core $\{0,1716,4862,16302,17160\}$, nine additional hard
residues $\{2431,9009,9867,17017,21164,24453,29315,31603,38896\}$
reopened by the exact-v2 recertification audit, and $27$ residues
whose closures under the Bridge-3 peeled-cofactor mechanism were
invalidated by the exact-v2 audit (and remain formally open pending
a sound replacement). The entire formal
chain is stated conditional on a single intentional axiom
\texttt{openResiduesStage1\_no\_solution} asserting that no $N$ with
$N\bmod M$ in the open set yields a \#647 candidate.

On each of the historical hard-$5$ classes, the divisor constraints
reduce the problem to simultaneous primality of an explicit
$13$-, $15$-, or $17$-term linear tuple in a parameter
$t$: a mod-$11$ square obstruction upgrades the budget-$\le 3$ forms
to primes, and residue-specific divisibility by $17$ or $19$ yields
additional prime-quotient constraints. We prove a completeness
theorem: at the budget bound $B_k\le 5$, the resulting list of
forced prime-or-prime-quotient conditions consists of exactly one
universal extension and six residue-specific triples, and this
bound is sharp.

We then prove that the operative local congruence program is
exhausted: no higher-power obstruction arises at $17$ or $19$, and
non-operative primes act only by residue translation. We introduce
the \emph{offset atlas}, an exact per-prime branch classification
of the affine forms $F_{r,k}(t)=2520(Mt+r)-k$; the atlas unifies
the mod-$11$ square-kill and the residue-specific extensions under
a single framework and provides machine-readable data against
which the divisor-sum conditions at every integer offset $k$ can
be tested. An empirical investigation of the
Erd\H{o}s-inequality-failure pattern yields \emph{killer sets}:
small sets $S_r\subseteq[1,100]$ such that every observed
strong near-miss at residue $r$ fails the Erd\H{o}s inequality
at some $k\in S_r$. Freezing the empirical claims to the audited
$301$-record strong $r=0$ corpus (and the combined $303$-record
strong corpus across $r\in\{0,1716,4862\}$), the size-$3$ set
$\{15,16,27\}$ is the unique minimum cover at $r=0$ ($301/301$),
while the empirical monitor $\{7,11,13,16\}$ covers the audited
strong corpus completely. A confirmed $13$-prime hit at
$t=11{,}741{,}230{,}040{,}310$ is consistent with
Bateman--Horn's leading-order quantitative prediction at the current
search range.

We further identify the structural mechanism governing the
empirical obstruction: on each atlas branch, the Erd\H{o}s
inequality at offset $k$ reduces to a divisor-count constraint
$\tau(Z_{r,k,\lambda}) \le B^{\mathrm{eff}}_{r,k,\lambda}$ on the
primitive cofactor, where the branchwise effective budget
$B^{\mathrm{eff}}_{r,k,\lambda} = \lfloor (k+2)/\tau(A_{r,k,\lambda})\rfloor$
is computable from the atlas alone. The branchwise severity
ordering induced by this criterion matches empirical pass-rate
ordering across all $87$ offsets $k\le 100$ at $r=0$ with Spearman
correlation $0.93$--$0.94$ under both empirical and neutral
branch weights, stable under chronological splitting
($\rho = 0.96$ between halves) and, for the neutral-prediction audit,
bootstrap resampling (median $\rho = 0.93$, $95\%$ CI
$[0.92,\,0.94]$). The monitor set $\{7,11,13,16\}$ has
an explicit structural interpretation: $k=16$ is atlas-tight at
every hard residue, while $k=7,11,13$ provide residue-specific
tightness at the residues where their primes are in the atlas
active set. A negative structural result is also reported: three
tested finite-incompatibility programs (covering at small primes,
pattern-rigidity, pairwise gcd identities) do not produce a
contradiction in the killer-set system at $r=0$. The remaining
gap is analytic: a density-level control on the budget-$\ge 4$
tail along Dickson sets. The nearest literature-shaped proxy
target is a tuple-restricted Titchmarsh-type divisor bound on
Dickson sets. At the reference hardest residue $r=0$ (tuple
rank~$13$) we establish the \emph{absolute form} of this proxy
(Remark~\ref{rem:track2_absolute}):
$\sum_{t\le X,\,t\in D_0(X)} \tau(L_{0,k_0}(t)) \ll X\,(\log X)^{-12}$
via a Selberg $\Lambda^{2}$ upper-bound sieve combined with
Shiu/Nair--Tenenbaum Brun--Titchmarsh. This absolute bound is
unconditional. The bridge from the absolute form to a density-level
statement inside the rank-$13$ Dickson set sits inside the same
Dickson/Schinzel-Hypothesis-H conjectural frame that
Proposition~\ref{prop:dichotomy} already uses to state the Track-$2$
dichotomy; it is flagged in Remark~\ref{rem:track2_absolute} as a
conditional outlook, not as a theorem of this paper. The analogous statements at
the other four hard residues are expected to hold
by the same composite but require residue-specific bookkeeping
(rank~$15$ or~$17$ sieve integrals, per-residue admissibility,
per-residue resultant support) which is not executed here.
A separate Dirichlet-$L$
pair-correlation route is secondary rather than load-bearing
here: in the current project notes it gives a conditional
asymptotic theorem under GRH plus a $\chi$-twisted
Hardy--Littlewood hypothesis $H_\chi$, with a finite-$T^*$
numerical corollary under an additional bookkeeping gate $B_\chi$;
the strict GRH-only finite-$T^*$ individual-$\chi$ version is
effectively red in the present framework. Thus the remaining
content is reduced to five explicit admissibility problems whose
small-offset obstruction profile is now structurally characterized
by the atlas's branchwise effective budget distribution.
\end{abstract}

\tableofcontents

%% ====================================================================
\section{Problem statement}\label{sec:problem}
%% ====================================================================

Let $\tau(m)$ denote the number of positive divisors of~$m$.

\begin{quote}
\textbf{Erd\H{o}s Problem \#647.} Does there exist $n>24$ such that
\[
  \max_{1\le m<n}(m+\tau(m))\le n+2\,?
\]
\end{quote}

Equivalently, $\tau(n-k)\le k+2$ for all $1\le k<n$. The known
solutions for $n\le 24$ are $\{1,2,3,4,5,6,8,10,12,24\}$; none
exist below $10^{10}$ (OEIS A087280). The problem remains open.

Erd\H{o}s conjectured the stronger statement
$\max_{m<n}(\tau(m)+m-n)\to\infty$ (recorded as
\texttt{variants.lim} in the DeepMind formal conjectures
repository~\cite{DeepMind}), which would imply at most finitely
many solutions.

A recent submission~\cite{Agbanwa} proposed a negative resolution via
a record-holder domination framework; the unresolved step is the
asymptotic claim that domination intervals eventually cover all
sufficiently large~$n$. The companion Lean artifact
axiomatized that asymptotic step rather than proving it, a
structural gap that Tao~\cite{TaoForum647} flagged as yielding
``no certificate of unconditional correctness.'' Our development
avoids this failure mode in a specific sense: the Lean chain contains a
single intentional axiom, \texttt{openResiduesStage1\_no\_solution},
whose scope is exactly the $41$ residues in the explicit finite set
\texttt{openResiduesStage1} (defined as the complement of the $55$
Lean-closed residues inside the $96$-class sieve frontier; see
Section~\ref{sec:formal}). That scope is computable and enumerable;
it is not a blanket ``all sufficiently large $n$'' axiom.
For four of the five historical hard residues
$\{1716, 4862, 16302, 17160\}$ (a subset of
\texttt{openResiduesStage1}), a $[1, 10^{15})$ exhaustive parity scan
at $t \equiv r \pmod{M}$ provides empirical stopping certificates
recording $0$ full-tuple witnesses.\footnote{Search methodology,
summarized from
\texttt{analysis/k12\_k13\_structure/findings/hard5\_stopping\_certificates\_20260423.md}:
per-residue candidate iteration under \texttt{--restrict-modM}
(only $t\equiv r\pmod{M}$), Eratosthenes sieve to limit $10^3$ with
primes dividing $M$ auto-skipped, BPSW primality pre-test on
sieve-survivors, then per-form deterministic Miller--Rabin with
$12$ bases (valid up to $3.3\times 10^{24}$) on each $\alpha_i t +
\beta_i$. Each residue processed $\approx 2.17\times 10^{10}$
candidate $t$-values; combined across the four residues the campaign
primality-tested on the order of $10^{5}$ sieve-survivors.}
We do not claim these stopping certificates are proofs; they bound
where counterexamples within those four classes could live and are
reported as empirical evidence supporting the axiom's scope.

Two elementary forum derivations (Dutta~\cite{Dutta} on 18 Jan 2026
and Alexeev~\cite{Alexeev} on 19 Jan 2026) independently arrive at
the same $n=2520N$ structure and the same tuple of forms
$d_iN-1$ that our Section~\ref{sec:shapes} derives in full (modulo
the subsequent refinements $315N-1\in\{p,2p\}$ and
$105N-1\in\{p,p^{2},2p,4p\}$ from Dutta's 19 Jan follow-up). Our
$n=2520N$ reduction and the forced-prime conditions at
$k=1,\ldots$ are consistent with theirs, and our further reduction
to $96$ classes modulo $M$ can be read as a formalization of
Alexeev's conjecture $11\mid n$, $13\mid n$ into the explicit
mod-$46189$ frame. The \emph{explicit near-witness}
$n = 604{,}517{,}614{,}941{,}240$ ($N = 239{,}887{,}942{,}437$)
exhibited by Alexeev satisfies $\tau(n-k)\le k+2$ for
$1\le k\le 13$ and $k=24$, demonstrating that the pointwise
tuple conditions alone are not contradictory; the analytic content
of \#647 is concentrated in the uniformity across all shifts $k<n$.

Our contribution is a structural reduction that localizes any
failure of asymptotic coverage to the $41$ residues
(decomposing as the historical hard-$14$ subset plus $27$ residues
whose prior closures were invalidated by the exact-v2 audit's
correction of the Bridge-3 peeled-cofactor mechanism).
Section~\ref{sec:hard} analyses the historical hard-$5$ structural
core in detail; Sections~\ref{sec:chain}--\ref{sec:formal} record
the formal chain.

%% ====================================================================
\section{Verified reduction}\label{sec:reduction}
%% ====================================================================

\begin{enumerate}
\item[\textbf{(R1)}] Every solution $n>84$ has $n=2520N$
  (Lean-verified: \texttt{erdos647\_div\_2520}).
\item[\textbf{(R2)}] Twelve shape constraints reduce
  $N\bmod{46189}$ to exactly $96$ residue classes
  (Lean-verified: \texttt{surviving\_card}).
\item[\textbf{(R3)}] Under exact-v2 closure semantics, adopted after
  the pointwise-unsound Bridge-3 peeled-cofactor transfer was retired
  (see Section~\ref{sec:formal} and the post-retirement audit at
  \texttt{docs/audits/55\_closed\_post\_bridge3\_audit\_20260424.md}),
  $55$ of $96$ classes are eliminated by Lean-checked closure
  certificates operating directly on the full value
  $m=2520\cdot N-k$:
  \begin{itemize}
  \item $42$ residues by \emph{single-overlap} arguments combining
    divisor monotonicity $\tau(d)\le\tau(m)$ when $d\mid m$ with
    sound coprime-multiplicativity $\tau(m)\ge 2\tau(a)$ under an
    explicit hypothesis $\gcd(p,a)=1$ (Lean-verified:
    \texttt{singleOverlapResidues42});
  \item $13$ residues by \emph{direct full-value} closure, the same
    full-value toolkit specialised to fixed content
    $g_{\mathrm{exact}}\in\{1729, 2584, 4522, 4845, 9690\}$
    (Lean-verified: \texttt{directFullValueResidues13}).
  \end{itemize}
  The remaining $41$ residues form the proof-critical open set
  \texttt{openResiduesStage1} $= \texttt{survivingResidues96}
  \setminus\texttt{closedResiduesStage1}$.
\item[\textbf{(R4)}] The $41$-residue open set decomposes as:
  \begin{itemize}
  \item the historical hard-$14$ core
    \[
      \{0,1716,2431,4862,9009,9867,16302,17017,17160,
        21164,24453,29315,31603,38896\}
    \]
    (the original hard-$5$ structural subset
    $\{0,1716,4862,16302,17160\}$ plus nine residues reopened by the
    exact-v2 recertification audit);
  \item the $27$ residues whose closures under the now-invalidated
    Bridge-3 peeled-cofactor mechanism require a sound replacement.
  \end{itemize}
  The Stage-1 axiom \texttt{openResiduesStage1\_no\_solution}
  asserts that no $N$ with $N\bmod M$ in this $41$-residue set
  yields a \#647 candidate.
\end{enumerate}

%% ====================================================================
\section{The 12-shape almost-prime package}\label{sec:shapes}
%% ====================================================================

\begin{proposition}[Extended shape package]
If $n=2520N$ satisfies the hypothesis, then twelve linear forms $dN-1$
($d\in\{105,140,210,252,280,315,420,504,630,840,1260,2520\}$) satisfy
explicit almost-prime shape conditions with bounded exponents.
\end{proposition}

The six pure-prime shapes ($k=1,2,3,4,6,12$) and six almost-prime
shapes ($k=5,8,9,10,18,24$) are derived by case analysis of
$\tau(2520N-k)\le k+2$. Shape elimination promotes $k=20$ from
budget-$3$ to budget-$2$, giving $13$ prime forms at $r=0$.
The full derivation is formalized in \texttt{Erdos647BridgeV1.lean}
(863~lines).

At each hard residue $r\in\{0,1716,4862,16302,17160\}$, the
budget-$\le 3$ prime tuple contains a universal $13$-form core
indexed by $k\in\{1,2,3,4,5,6,8,9,10,12,18,20,24\}$ with
coefficients $\alpha_i=d_i\cdot 46189$ for $d_i\in\{2520,1260,
840,630,504,420,315,280,252,210,140,126,105\}$. Residue-specific
extras extend the full tuple to sizes $13,17,15,15,17$ for
$r\in\{0,1716,4862,16302,17160\}$ respectively; these extras
arise when $\gcd(\alpha_k,\beta_k^{(r)})$ picks up residue-dependent
factors of the sieve primes $\{11,13,17,19\}$ absent from the
core. See \texttt{data/hard\_root\_forms\_r*.json} for the
per-residue breakdown.

%% ====================================================================
\section{Closure mechanisms}\label{sec:closure}
%% ====================================================================

\subsection{Budget-$\le 1$ contradiction}

For $72$ residues, some $k$-value produces a form whose total content
$g=\gcd(2520\cdot 46189,\,|2520r-k|)$ satisfies $\tau(g)>k+2$,
making the constraint $\tau(2520N-k)\le k+2$ impossible.

\subsection{Budget-$2$ inadmissibility}

For $5$ residues, the set of budget-$2$ forms (those requiring
$\tau(\text{quotient})\le 2$, i.e., primality) covers all of
$\mathbb{Z}/p\mathbb{Z}$ at some prime~$p$. Every $t$ has at
least one budget-$2$ form divisible by~$p$, hence composite.

\subsection{Budget-$2$/$3$ inadmissibility (b23 mechanism)}

\begin{theorem}[b23 branch-closure, Lean-verified]
Let $\{L_i(u)=a_i u+b_i\}$ be affine forms with budgets $2$ or~$3$.
If a prime~$p$ covers all residues mod~$p$ via these forms, and every
covering form satisfies $L_i(u)>p^2$ for $u\ge U$, then no $u\ge U$
satisfies all budget constraints simultaneously.
\end{theorem}

\begin{proof}
Fix $u\ge U$. Some covering form $L_i$ has $p\mid L_i(u)$.
If $\text{budget}=2$: $L_i(u)>p$ forces $\tau(L_i(u))\ge 3>2$.
If $\text{budget}=3$: $L_i(u)>p^2$ forces $\tau(L_i(u))\ge 4>3$
(since $1,p,L_i(u)/p,L_i(u)$ are four distinct divisors when
$L_i(u)/p>p$). Either contradicts the budget.
\end{proof}

For $10$ residues, the b23 mechanism applies at the $t$-level
(without branching). For $4$ residues ($r\in\{858,2574,24310,25168\}$),
branching $t=30u+c$ followed by b23 at the $u$-level closes all
$30$ branches.

%% ====================================================================
\section{Structural characterization of the hard set}\label{sec:hard}
%% ====================================================================

\begin{proposition}[Necessary condition, Lean-verified]
Every hard residue satisfies $143\mid r$, where $143=11\cdot 13$.
\end{proposition}

This has a clean structural consequence: since $11\cdot 13\mid r$
and $7$ divides every $d$-coefficient, the forms
$\alpha_i=d_i\cdot 46189$ satisfy $\alpha_i\equiv 0\pmod{q}$ for
$q\in\{7,11,13\}$. Hence these three sieve primes contribute
\emph{zero} forbidden residues, nullifying three of four operative
sieve primes. Only $\{2,3,5,17,19\}$ remain.

\begin{remark}
The condition $143\mid r$ is necessary but not sufficient: $43$ of
$48$ survivors with $143\mid r$ close via budget or inadmissibility
mechanisms at the remaining primes. The $5$ hard residues are
those where even primes $\{2,3,5\}$ fail to produce a covering
or budget contradiction.
\end{remark}

\begin{remark}[Sieve affinity for $143$-divisibility]\label{rmk:sieve-affinity}
Of the $96$ surviving residues mod $46189$, exactly $48$ (that is,
$50\%$) are divisible by $143$. In the ambient space mod $46189$,
only $323/46189\approx 0.70\%$ of residues are $143$-divisible.
The $12$-shape sieve thus exhibits a $71.5\times$ affinity for
$143$-divisible residues, consistent with the observation that
$143$-divisibility nullifies the sieve's most aggressive primes
($11$ and $13$). Hardness, once $143$-divisibility is present,
is then decided among these $48$ residues by the behaviour at the
residual operative primes $\{2,3,5,17,19\}$.
\end{remark}

\begin{proposition}[Operative-prime staticity at $17$ and $19$]\label{prop:staticity}
For every shape $k$ in the $13$-form core tuple and every residue
$r$ mod $46189$, the leading coefficient $\alpha_k = d_k\cdot 46189$
satisfies $\alpha_k\equiv 0\pmod{17}$ and $\alpha_k\equiv 0\pmod{19}$.
Consequently
\[
  L_{r,k}(t)\equiv\beta_k^{(r)}\pmod{17}
  \qquad\text{and}\qquad
  L_{r,k}(t)\equiv\beta_k^{(r)}\pmod{19}
\]
for every $t\in\mathbb{Z}$: divisibility of $L_{r,k}(t)$ by $17$ or
$19$ is determined entirely by the constant term
$\beta_k^{(r)}$, independently of the $t$-parameter.
\end{proposition}

Since $46189 = 11\cdot 13\cdot 17\cdot 19$ and every core
$d$-coefficient is a multiple of $7$, the factorisation
$\alpha_k = d_k\cdot 46189$ forces $\alpha_k\equiv 0$ modulo each of
$\{7,17,19\}$ as a structural feature of the reduction itself,
independent of $r$. The $143$-divisibility condition of the
necessary-condition proposition adds $\{11,13\}$ to this list
whenever $143\mid r$. Among the five operative sieve primes
$\{2,3,5,17,19\}$ that remain after nullifying $\{7,11,13\}$,
Proposition~\ref{prop:staticity} thus isolates $\{17,19\}$ as
carrying no $t$-parameter information: their contribution to
closure or hardness is determined \emph{rigidly} by the constants
$\beta_k^{(r)}\bmod 17$ and $\beta_k^{(r)}\bmod 19$. This is the
structural content that the CRT factorization
(Proposition~\ref{prop:crt}) observes.

\begin{proposition}[CRT factorization of hardness, Lean-verified]\label{prop:crt}
Let $\pi\colon \mathbb{Z}/46189\mathbb{Z}\to
(\mathbb{Z}/17\mathbb{Z})\times(\mathbb{Z}/19\mathbb{Z})$ be the
natural CRT projection. Restricted to the $48$ surviving residues
with $143\mid r$, the projection $\pi$ is injective on both the
hard and closing subsets, and their images are disjoint:
\[
  \pi(\mathrm{hardResidues5}) \;=\; \mathcal{H}_{17,19} \;:=\;
  \{(0,0),\,(16,6),\,(0,17),\,(16,0),\,(7,3)\},
\]
\[
  \pi(\mathrm{closingResidues}) \cap \mathcal{H}_{17,19} \;=\; \emptyset.
\]
Consequently, for $r\in\mathrm{survivingResidues}$,
\[
  r\in\mathrm{hardResidues5}
  \;\iff\;
  143\mid r \;\text{and}\; \pi(r)\in\mathcal{H}_{17,19}.
\]
(Lean theorem \texttt{hard\_residue\_crt\_factorization},
\texttt{native\_decide}.)
\end{proposition}

The content of Proposition~\ref{prop:crt} is the injectivity and
disjointness, not the enumeration of $\mathcal{H}_{17,19}$: the
statement is that the $(\bmod\,17,\bmod\,19)$ projection is
sufficient statistics for hardness, given $143$-divisibility. Any
future structural characterization of the historical hard-$5$ subset within the original $91/5$
partition must factor through this projection. The hardness of the
five residues therefore appears to arise from the simultaneous
failure of all five operative closure mechanisms at
$\{2,3,5,17,19\}$, rather than from a unifying structural property.
Whether $\mathcal{H}_{17,19}$ admits a deeper arithmetic
characterization is an open question we flag for the community.

\begin{remark}[Axis overrepresentation in $\mathcal{H}_{17,19}$]\label{rmk:axes}
Within $\mathcal{H}_{17,19}$, three of the five pairs lie on a
coordinate axis — $(0,0)$, $(0,17)$, and $(16,0)$ — compared to
only $17+19-1=35$ of $17\cdot 19 = 323$ pairs (roughly $10.8\%$)
in the ambient space. The hard set is therefore about $5.5\times$
overrepresented on the axes. Arithmetically, this says three of
the five hard residues have an \emph{additional} divisibility by
one of the residual sieve primes $17$ or $19$ on top of
$143$-divisibility: $r=0$ is divisible by both $17$ and $19$,
$r=4862$ is divisible by $17$, and $r=16302$ is divisible by $19$.
The remaining two hard residues, $r=1716$ and $r=17160$, have
neither additional divisibility but nonetheless project into
$\mathcal{H}_{17,19}$ for reasons we have not characterized
structurally.
\end{remark}

\begin{remark}[Jacobi-sign necessary condition]\label{rmk:jacobi}
Complementing the axis pattern, a cross-Jacobi observation:
for each of the $48$ surviving residues $r$ with $143\mid r$,
form the pair of Jacobi symbols
$\bigl(\tfrac{r}{17},\,\tfrac{r}{19}\bigr)$ and their product
$\bigl(\tfrac{r}{17}\bigr)\cdot\bigl(\tfrac{r}{19}\bigr)\in\{-1,0,+1\}$.
Of the $48$ surviving $143$-divisible residues, exactly $18$ have
product $-1$ (the ``mixed QR/NR'' case, where one of
$r\bmod 17, r\bmod 19$ is a quadratic residue and the other is a
non-residue), and \emph{zero} of those $18$ are hard; the five hard
residues all have product in $\{0, +1\}$. Explicitly:
\[
\Bigl(\tfrac{r}{17}\Bigr)\cdot\Bigl(\tfrac{r}{19}\Bigr)
\;=\;\left\{
\begin{array}{ll}
0,     & r\in\{0,\,4862,\,16302\},\\
+1,    & r\in\{1716,\,17160\}.
\end{array}
\right.
\]
Hence a partial necessary condition: \emph{if
$r\in\mathrm{survivingResidues}$, $143\mid r$, and the Jacobi
product equals $-1$, then $r$ is closing}. The condition is not
sufficient — $15$ of the $30$ surviving $143$-divisible residues
with non-mixed Jacobi product are also closing — but it strictly
reduces the search space for hardness by $60\%$ (from $48$ to $30$
candidates). The observation is presently empirical; Lean
formalization would require routine \texttt{native\_decide}-style
enumeration. Data: \texttt{data/hard\_residue\_characterization.json}.
\end{remark}

\begin{proposition}[Branch-level hard set, Lean-verifiable]\label{prop:branch18}
Among the $150$ nominal level-$1$ branch-states $(r,c)$ with
$r\in\mathrm{hardResidues5}$ and $c\in\mathrm{Fin}\ 30$, exactly
$18$ pairs genuinely resist closure at budget $\le 3$:
\begin{align*}
S_{18} \;=\; \{&(0,0),\,(0,8),\,(0,10),\,(0,18),\,(0,20),\,(0,28),\\
                &(1716,6),\,(1716,16),\\
                &(4862,12),\,(4862,22),\\
                &(16302,0),\,(16302,2),\,(16302,10),\,(16302,12),\,(16302,20),\,(16302,22),\\
                &(17160,0),\,(17160,18)\}.
\end{align*}
The remaining $132$ pairs admit a budget-$\le 3$ inadmissibility
certificate covering some operative prime (budget-$\le 1$ directly,
or covering of $\mathbb{Z}/p\mathbb{Z}$ at $p\in\{2,3,5\}$ by
budget-$\le 2$ or budget-$\le 3$ forms). The cardinality claim
$|S_{18}| = 18$ and the residue-distribution claim
$6 + 2 + 2 + 6 + 2$ are Lean-verifiable
(\texttt{hardBranches18\_card},
\texttt{hardBranches18\_distribution} in
\texttt{Erdos647ReductionChain.lean}, both by \texttt{decide}).
\end{proposition}

\begin{remark}[Distribution asymmetry]\label{rmk:asym}
The $18$ tentative hard branches distribute asymmetrically across
the $5$ hard residues: $6$ each at $r=0$ and $r=16302$, and $2$ each
at $r=1716$, $r=4862$, $r=17160$. The tentative-covering prime and
required budget cap break down as: $11$ branches at $(p=2$,
budget $\le 6)$; $5$ branches at $(p=3$, budget $\le 4)$; $2$
branches at $(p=2$, budget $\le 4)$. Whether the
$r=0 / r=16302$ concentration reflects a deeper arithmetic feature
distinguishing these two residues from the other three hard
residues is an open question. Enumeration computed in
\texttt{scripts/higher\_budget\_enumeration.py}; data recorded in
\texttt{data/tentative\_hard\_branches.json}.
\end{remark}

\begin{remark}[Scope of pre-analytic methods]\label{rmk:pre-analytic}
We tested three natural refinements that a sieve-theoretic attack
might attempt against the $18$ tentative branches, each designed to
extract additional structure from the reduction without introducing
new analytic hypotheses. Each is ruled out; together they
characterize the boundary of pre-analytic methods on this set.

\emph{(i) Finer branching at primes beyond $\bmod 30$.} Moduli
$M_{\mathrm{new}} \in \{210, 510, 570, 714, 798, 1938\}$
(incorporating primes $7$, $17$, $19$ in various combinations) were
tested on each of the $18$ tentative branches. Under no such
$M_{\mathrm{new}}$ do all $M_{\mathrm{new}}/\gcd(M_{\mathrm{new}},30)$
sub-branches close rigorously at budget $\le 3$. Richer branching
produces partial closures — for example, $(r, c) = (17160, 0)$ under
$M_{\mathrm{new}} = 1938$ closes $18$ of its $64$ sub-branches
rigorously, but the remaining $46$ stay tentative. Consistent with
the barrier theorem (Section~\ref{sec:barrier}): no finite
prime-by-prime descent at primes greater than the clause count
can eliminate a tentative branch.

\emph{(ii) Shape restriction via square-elimination of budget-$3$
forms.} A budget-$3$ form $L(u) = \alpha u + \beta$ admits two
arithmetic shapes: $L = p$ (prime) or $L = p^2$ (prime square). The
$p^2$ option can be ruled out uniformly in $u$ whenever there is a
prime $q$ with $q \mid \alpha$, $q \nmid \beta$, and
$\beta \bmod q$ a non-zero quadratic non-residue modulo $q$; then
$L(u) \equiv \beta \pmod q$ is constant and a non-residue, hence
never a square. For every budget-$3$ form in every one of the $18$
tentative branches such a witness prime $q$ exists ($8/8$ through
$12/12$ forms eliminated per branch, $100\%$ elimination rate in all
$18$; see
\texttt{scripts/qnr\_elimination\_prototype.py}). The effective form
pool is thus reduced to budget $\le 2$ throughout; yet no branch
closes, because the budget-$\le 2$ forms still fail to cover
$\mathbb{Z}/p\mathbb{Z}$ at any operative prime in $\{2, 3, 5\}$.
This rules out the entire family of square-based enhancements —
QNR, Legendre symbol constraints, discriminant obstructions, any
variant preventing $L$ from being $p^2$ — because even the strongest
possible outcome of such an analysis leaves the budget-$\le 2$
coverage deficit unchanged. The closure obstruction is localized to
the budget-$2$ arithmetic itself.

\emph{(iii) Sub-branching on the covering prime itself.} For each
tentative branch with tentative covering prime
$p \in \{2, 3\}$, sub-branching mod $30p$ (mod $60$ for $p = 2$,
mod $90$ for $p = 3$) causes natural promotion of some forms to
effective budget $1$ under the reduction pipeline, yielding direct
contradictions in certain sub-branches. At $p = 3$ exactly $1$ of
the $3$ sub-branches closes this way per parent; at $p = 2$, $0$ of
the $2$ sub-branches close. Since parent closure requires all
sub-branches to close, no parent closes
(\texttt{scripts/subbranch\_promotion\_prototype.py},
data in \texttt{data/subbranch\_promotion\_attempts.json}). This
rules out the most natural "branch the deficit" attack as well.

\emph{(iv) Residue-level closure rates under operative versus
non-operative prime branching.} A finer-grained extension of~(i).
At $M_\mathrm{new} = 30 \cdot 17 \cdot 19 = 9690$ (adding both
operative primes dividing $M$), $488$ of the $5{,}814$ level-$2$
sub-branches close ($8.39\%$), with a sharp residue-dependent
hierarchy: $r=0$ at $0.00\%$, $r=16302$ at $5.88\%$, $r=1716$ at
$10.84\%$, $r=4862$ at $18.73\%$, $r=17160$ at $28.33\%$. The
immunity of $r=0$ has an arithmetic explanation: $2520r - k = -k$
forces $\gcd(2520M, |2520r-k|) = \gcd(2520, k)$ (every $k$ in the
$12$-shape set is at most $24$, hence coprime to
$M = 11\cdot 13\cdot 17\cdot 19$), so the base gcd contains no
factor of~$M$, $a_{\mathrm{orig}} = 2520M/g$ absorbs all of $M$, and
the extra-gcd peeling step has no $M$-prime factors to absorb under
refinement at primes dividing~$M$. Replacing the operative pair
$\{17,19\}$ with non-operative $\{23,29\}$ bypasses the obstruction:
at $M_\mathrm{new} = 30\cdot 23\cdot 29 = 20010$, $10{,}827$ of
$12{,}006$ sub-branches close ($90.18\%$), and the per-residue
hierarchy collapses to a tight $87.96\%$–$93.40\%$ band across all
five hard residues (in particular $r=0$ rises from $0\%$ to
$87.96\%$). The finding generalizes to other non-operative pairs: at
$M_\mathrm{new} = 30\cdot 31\cdot 37 = 34{,}410$, $16{,}526$ of
$20{,}646$ sub-branches close ($80.04\%$) with the per-residue range
$75.62\%$–$83.70\%$. Smaller non-operative primes are more effective
(each prime gives richer $\gcd(a_u, b_u)$ structure per unit $j$), but
the mechanism is not specific to $\{23, 29\}$.

This shows that the residue-level variation observable at mod-$9690$
is a reflection of $M$-factor availability in the reduction pipeline,
not an intrinsic hardness hierarchy among the five hard residues; the
$(r \bmod 17, r\bmod 19)$ pair structure from
Proposition~\ref{prop:crt} predicts the mod-$9690$
response but not the mod-$20010$ or mod-$34410$ response.

\emph{Residual structure at $M_\mathrm{new}{=}20010$.} The $1{,}179$
open sub-branches left by non-operative branching have a sharp
mod-$23$ signature: of the $23$ residue classes $j \bmod 23$ (where
$j$ is the level-$2$ branching index), the class $j \equiv 8 \pmod{23}$
has $65.5\%$ of its members universally-closing (closing at every one
of the $18$ tentative $(r,c)$), while $j \equiv 7 \pmod{23}$ has $0\%$
and is the unique class with no universally-closing member. The
mechanism: $30 \equiv 7 \pmod{23}$, so the $23$-factor peeling in
$\gcd(a_u, b_u) = \gcd(20010\,a_{\mathrm{orig}},
a_{\mathrm{orig}}(30 j + c) + b_{\mathrm{orig}})$ fires iff $j$ lies
in a specific class mod~$23$ determined by $(r, c, k)$; the class
$j \equiv 7 \pmod{23}$ is precisely the one receiving no hit across
the twelve shapes. The residual hardness of non-operative branching
is thus a single CRT bottleneck, paralleling the $r=0$ immunity at
mod-$9690$ but at the level-$2$ index rather than the residue.

Consistent with~(i), no single $M_\mathrm{new}$ closes every
sub-branch of any tentative~$(r,c)$, so no full tentative branch is
eliminated by branching alone; (i)'s conclusion — that no finite
prime-by-prime descent at primes greater than the clause count can
eliminate a tentative branch — is reaffirmed. Experimental details:
\texttt{scripts/a1\_branch\_323.py},
\texttt{scripts/a1b\_branch\_667.py},
\texttt{scripts/a1d\_branch\_1147.py},
\texttt{scripts/a1b\_residual\_structure.py};
memos \texttt{docs/a1\_branch323\_result.md},
\texttt{docs/a1b\_branch667\_result.md},
\texttt{docs/c1\_pair\_partition.md},
\texttt{docs/a1b\_residual\_structure.md}.

Together, (i)–(iv) exhaust the natural pre-analytic routes to
closing the $18$ tentative branches. The closure obstruction
localizes to the budget-$\le 2$ coverage at operative primes
$\{2, 3, 5\}$, which no combination of finer branching (at operative
or non-operative primes), form-shape restriction, or sub-branching
on the covering prime can affect.
Further progress on this set appears to require analytic input,
specifically results of the form discussed in
Section~\ref{sec:conditional} — a density-level control on the
budget-$\ge 4$ Dickson-set tail, with a tuple-uniform
Titchmarsh-type divisor bound as the nearest proxy target, or a
prime-core-relative analog of Matthiesen's affine-linear
correlation estimates on the $13$-form Dickson set.
\end{remark}

%% ====================================================================
\section{The mod-$11$ square-kill and divisor-extension theorem}\label{sec:div_ext}
%% ====================================================================

The hard-set characterization of Section~\ref{sec:hard} leaves a
congruential residue of open content: for each hard
$r\in\{0,1716,4862,16302,17160\}$ we know the $12$-shape package is
budget-$\le 3$-admissible on the progression $N=Mt+r$, but the package
is not pinned to primality until we close the mod-$11$ shape
classification. This section supplies that pinning via a standalone
arithmetic lemma on integers $x\equiv -1\pmod{11}$, derives a
divisor-extension corollary that applies to Erd\H{o}s $\#647$
candidates on hard residues, and proves a completeness theorem:
under the budget bound $B_k\le 5$, the resulting list of forced
primality constraints consists of exactly one universal extension
and six residue-specific quotient-prime extensions.

\subsection{The mod-$11$ arithmetic lemma}

Throughout, write $M=46189=11\cdot 13\cdot 17\cdot 19$ and, for
$k\mid 2520$, set $d_k := 2520/k$ and
\[
  x_k(N) \;:=\; d_k\cdot N - 1,
  \qquad
  B_k \;:=\; \bigl\lfloor (k+2)/\tau(k) \bigr\rfloor.
\]

\begin{lemma}[Mod-$11$ rigidity]\label{lem:mod11}
Let $x>1$ be an integer with $x\equiv -1\pmod{11}$.
\begin{enumerate}
\item If $\tau(x)\le 3$, then $x$ is prime.
\item If $q\in\{17,19\}$, $q\mid x$, and $\tau(x)\le 5$, then
$v_q(x)=1$ and $x/q$ is prime.
\end{enumerate}
\end{lemma}

\begin{proof}
(i) An integer with $\tau\le 3$ has shape in $\{1,p,p^2\}$. The
quadratic residues mod~$11$ are $\{1,3,4,5,9\}$; the class
$-1\equiv 10\pmod{11}$ is a non-residue. Hence $x$ is never a
perfect square, ruling out $p^2$, while $x>1$ excludes $1$. So
$x=p$ is prime.

(ii) An integer with $\tau\le 5$ divisible by $q\in\{17,19\}$ has
shape $q^e$ for some $1\le e\le 4$, or $q\cdot p$ with $p$ prime and
$p\ne q$. Direct computation gives
\[
  17^e \pmod{11}:\;\; 6,\,3,\,7,\,9 \qquad (e=1,2,3,4),
\]
\[
  19^e \pmod{11}:\;\; 8,\,9,\,6,\,4 \qquad (e=1,2,3,4);
\]
none of these equals $10\equiv -1$. The shapes $q^e$ ($1\le e\le 4$)
are therefore excluded by $x\equiv -1\pmod{11}$. The only remaining
possibility is $x=q\cdot p$, whence $v_q(x)=1$ and $x/q=p$ is prime.
\end{proof}

\subsection{Divisor-extension on Erd\H{o}s candidates}

Write $K_{\text{pure}}:=\{1,2,3,4,6,12\}$ for the divisors of $2520$
with $v_p(k)<v_p(2520)$ for every prime $p\mid k$.

\begin{corollary}[Divisor-extension at hard residues]\label{cor:div_ext}
Fix a hard residue
\[
  r\in\{0,1716,4862,16302,17160\},
\]
and let $N=Mt+r$ for some $t\ge 0$. Suppose $n=2520N$ is an
Erd\H{o}s $\#647$ candidate.
\begin{enumerate}
\item For $k\in K_{\text{pure}}$, $\gcd(k,x_k(N))=1$
unconditionally; hence $\tau(n-k)=\tau(k)\cdot\tau(x_k(N))$ and the
Erd\H{o}s inequality at $k$ rewrites as $\tau(x_k(N))\le B_k$.
\item $x_k(N)\equiv -1\pmod{11}$ for every $k\mid 2520$.
\item For $k\in K_{\text{pure}}$ with $B_k\le 3$, $x_k(N)$ is
prime. These are the six core budget-$\le 3$ forms with
$d\in\{2520,1260,840,630,420,210\}$.
\item For $k\not\in K_{\text{pure}}$, coprimality can fail and the
budget constraint combined with the per-prime case analysis of
Section~\ref{sec:shapes} gives the shape conditions recorded there.
\item For $k=20$ ($d_{20}=126=2\cdot 3^2\cdot 7$), coprimality
$\gcd(20,x_{20}(N))=1$ holds except on the sub-progression
$N\equiv 1\pmod 5$ (Case~B), where $5\mid x_{20}(N)$. Case~A
($N\not\equiv 1\pmod 5$) forces $x_{20}(N)$ prime; Case~B forces
$x_{20}(N)=5\cdot q$ with $q$ prime and $q\equiv 2\pmod{11}$.
\end{enumerate}
\end{corollary}

\begin{remark}[Erd\H{o}s candidacy is a genuine hypothesis]
The candidate assumption in Corollary~\ref{cor:div_ext} is necessary.
At $N=M=46189$ (a hard-residue survivor with $r=0$ that is not an
Erd\H{o}s candidate) the form
$x_2(N)=1260\cdot 46189-1=58{,}198{,}139=5297\cdot 10987$
has $\tau(x_2)=4>2=B_2$. Survivorship alone gives the mod-$11$
congruence but not the $\tau$-bound; candidacy supplies the bound.
\end{remark}

\subsection{Tightness of the budget bound}

\begin{proposition}[Tightness of $B_k\le 5$]\label{prop:tight}
The budget bound $B_k\le 5$ in Corollary~\ref{cor:div_ext}(v) and
in Theorem~\ref{thm:complete} below is sharp: there exist
hard-residue integers $N$ and divisors $k\mid 2520$
with $B_k=6$ for which $x_k(N)$ is divisible by $q\in\{17,19\}$
yet $x_k(N)/q$ is composite, while the local Erd\H{o}s inequality at
that single offset $k$ is still satisfied.
\end{proposition}

\begin{proof}
Take $r=1716$, $k=72$, $d_k=35$, $B_{72}=\lfloor 74/\tau(72)\rfloor=\lfloor 74/12\rfloor=6$, and $t=42$. Then
$N=Mt+r=46189\cdot 42+1716=1{,}941{,}654$, and
$x_{72}(N)=35N-1=67{,}957{,}889=19^2\cdot 188{,}249$ with
$188{,}249$ prime. Thus $\tau(x_{72})=6=B_{72}$, $x_{72}\equiv 10\pmod{11}$, and the Erd\H{o}s inequality at $k=72$ is satisfied
($\tau(n-72)=\tau(72)\tau(x_{72})=12\cdot 6=72\le 74=k+2$). Yet
$v_{19}(x_{72})=2$, so $x_{72}/19=19\cdot 188{,}249$ is composite,
violating the conclusion of Corollary~\ref{cor:div_ext} if it were
extended to $B_k\le 6$.

The failure is structural. At $\tau(x)=6$ two new shapes appear
beyond those admitted at $\tau(x)\le 5$: $p^5$ and $p^2q$. (The
shape $pqr$ has $\tau=8$ and is not in play.) Both new shapes are
compatible with $x\equiv -1\pmod{11}$: direct computation yields
$17^5\equiv 19^5\equiv -1\pmod{11}$ (admitting $p^5$ for
$p\in\{17,19\}$), and $q\cdot s^2\equiv -1\pmod{11}$ has solutions
in $s$ for each $q\in\{17,19\}$. So mod-$11$ no longer forces the
prime-quotient conclusion at $B_k=6$.
\end{proof}

\subsection{Completeness at $B_k\le 5$}

\begin{theorem}[Completeness of the divisor-extension list]\label{thm:complete}
Among the $48$ divisors $d\mid 2520$:
\begin{enumerate}
\item Exactly thirteen divisors $d$ satisfy $B_{2520/d}\le 3$:
\[
  d\in
  \{105,140,210,252,280,315,420,504,630,840,1260,2520\}
  \cup\{126\}.
\]
The first twelve are the $12$-shape sieve set (Section~\ref{sec:shapes});
$d=126$ (from $k=20$, $B_{20}=3$) is the universal extension.
\item Exactly ten divisors $d$ satisfy $B_{2520/d}\in\{4,5\}$:
$d\in\{42,60,63,70,84,90,120,168,180,360\}$.
Among all pairs $(r,d)$ with $r$ hard and $d$ in this list, the
congruence $d\cdot r\equiv 1\pmod q$ (for $q\in\{17,19\}$) holds
in exactly the following six cases:
\[
  \renewcommand{\arraystretch}{1.15}
  \begin{array}{l|l|l|c}
    r & d & k & \text{forced by}\\ \hline
    1716  & 84  & 30 & 17 \\
    1716  & 168 & 15 & 19 \\
    4862  & 180 & 14 & 19 \\
    16302 & 84  & 30 & 17 \\
    17160 & 70  & 36 & 19 \\
    17160 & 90  & 28 & 17 \\
  \end{array}
\]
For each such $(r,d,q)$ and any hard-residue Erd\H{o}s candidate $N$
with $N\equiv r\pmod M$, the form $x_k(N)=d\cdot N-1$ is divisible
by $q$ and the quotient $x_k(N)/q$ is prime (possibly modulo a
Case-B-type sub-progression analog of Corollary~\ref{cor:div_ext}(v)
at a prime other than $q$).
\item No divisor-extension constraint arises for any $(r,d)$ outside
the thirteen universal and six residue-specific cases above.
\end{enumerate}
\end{theorem}

\begin{proof}
Part (i) is arithmetic: the thirteen divisors listed are those with
$B_{2520/d}\le 3$. Part (ii) enumerates the divisors $d\mid 2520$
with $B_{2520/d}\in\{4,5\}$ and tests $dr\equiv 1\pmod{17}$ and
$dr\equiv 1\pmod{19}$ across the $50$ candidate pairs
$(r,d)\in\mathrm{hardResidues5}\times\{42,60,\ldots,360\}$; the six
cases tabulated are the complete list. Part (iii) follows from the
exhaustion of $B_k\le 5$: at $B_k\ge 6$ the conclusion can fail
(Proposition~\ref{prop:tight}), and divisors with $B_{2520/d}\le 5$
are covered by (i) and (ii). The computational content is
enumerated in \texttt{scripts/verify\_omitted\_forms.py}.
\end{proof}

\begin{remark}[BH-friendly formulation]\label{rem:bhfriendly}
Expanding $N=Mt+r$ in each forced-prime form and clearing the
denominator where applicable, the Case~A pure-primality constraints
of Theorem~\ref{thm:complete} can each be written as ``$L(t)$ prime''
for an explicit integer-valued linear form $L\in\mathbb{Z}[t]$.
On the Case~A sub-progression (complementary to $N\equiv 1\pmod 5$
at $r=0$, and analogously elsewhere), the Erd\H{o}s $\#647$
candidate condition reduces on the hard residues to simultaneous
primality of an explicit tuple of linear forms in $t$: $13$ forms
at $r=0$, $14$ at $r=4862$, $14$ at $r=16302$, $15$ at $r=1716$,
and $15$ at $r=17160$. This is the Bateman--Horn-friendly
formulation used in Section~\ref{sec:BH}.
\end{remark}

%% ====================================================================
\section{The offset atlas}\label{sec:atlas}
%% ====================================================================

The Case-A/Case-B structure of Corollary~\ref{cor:div_ext} and the
residue-specific extensions of Theorem~\ref{thm:complete} become
transparent under a single coordinatization of the affine forms
$F_{r,k}(t):=2520(Mt+r)-k$, which we call the \emph{offset atlas}.

\begin{definition}[Atlas objects]
Write $A:=2520\cdot M$, $C_{r,k}:=2520\cdot r-k$, $P_0:=\{2,3,5,7,11,13,17,19\}$,
and $a_p:=v_p(A)$, $c_p(r,k):=v_p(C_{r,k})$. A prime $p\in P_0$ is
\emph{active} at $(r,k)$ iff $c_p\ge a_p$; otherwise \emph{passive}.
Write $S(r,k)$ for the set of active primes.
\end{definition}

For each active prime $p\in S(r,k)$, Hensel lifting produces a
unique residue class $t_{p,j}\pmod{p^j}$ at each level $j\ge 1$ on
which $v_p(F_{r,k}(t))\ge a_p+j$. A \emph{branch} is an exact-valuation
specification $\lambda=(\lambda_p)_{p\in S(r,k)}$: at $\lambda_p=j$,
we require $t\equiv t_{p,j}\pmod{p^j}$ and $t\not\equiv t_{p,j+1}\pmod{p^{j+1}}$
(so $v_p(F_{r,k}(t))=a_p+\lambda_p$ exactly). Combining across active
primes via CRT pins $t$ modulo $m_\lambda:=\prod_{p\in S(r,k)}p^{\lambda_p}$.
After reparameterizing $t=t_\lambda+m_\lambda u$, we obtain the
\emph{primitive cofactor} $Z_{r,k,\lambda}(u)=\alpha u+\beta\in\mathbb{Z}[u]$
with $\gcd(\alpha\cdot u+\beta,A)=1$ for $u$ in the branch.

\begin{proposition}[Atlas factorization]\label{prop:atlas_factorization}
On branch $\lambda$ of $(r,k)$,
\[
  F_{r,k}(t)
  \;=\;
  \Bigl(\prod_{p\notin S}p^{c_p}\Bigr)
  \cdot
  \Bigl(\prod_{p\in S}p^{a_p+\lambda_p}\Bigr)
  \cdot
  Z_{r,k,\lambda}(u),
\]
and
\[
  \tau(F_{r,k,\lambda}(t))
  \;=\;
  T^{\mathrm{fix}}_{r,k}\cdot\Bigl(\prod_{p\in S}(a_p+\lambda_p+1)\Bigr)\cdot\tau(Z_{r,k,\lambda}(u)),
\]
where $T^{\mathrm{fix}}_{r,k}:=\prod_{p\notin S}(c_p+1)$. Consequently
the branch budget
\[
  B_{r,k,\lambda}
  \;:=\;
  \Bigl\lfloor (k+2)\bigl/\bigl(T^{\mathrm{fix}}_{r,k}\prod_{p\in S}(a_p+\lambda_p+1)\bigr)\Bigr\rfloor
\]
bounds $\tau(Z_{r,k,\lambda}(u))$ whenever $N$ is an Erd\H{o}s
$\#647$ candidate.
\end{proposition}

The atlas is finite: restricting to $(r,k,\lambda)$ with
$B_{r,k,\lambda}\ge 2$, for $r$ hard and $k\in[1,100]$, produces
$500$ $(r,k)$ entries with a total of $3{,}471$ branches. At each
branch the atlas also records the residues $Z_{r,k,\lambda}(u)\pmod p$
for $p\in P_0$ (fixed when $\alpha\equiv 0\pmod p$, otherwise
varying with $u$), and the list of exponent-patterns of
$Z_{r,k,\lambda}(u)$ compatible with these residues via the
order-subgroup realizability criterion in $(\mathbb{Z}/p)^{\times}$.
The machine-readable dataset is in \texttt{data/offset\_atlas.json}.

\begin{remark}[Case~A vs Case~B at $k=20$]
At $(r,k)=(0,20)$, Proposition~\ref{prop:atlas_factorization} yields
$S=\{5\}$, $T^{\mathrm{fix}}=3$, and two admissible branches with
$\lambda_5\in\{0,1\}$. Case~A ($\lambda_5=0$) has budget~$3$ and
unique surviving pattern $[1]$; Case~B ($\lambda_5=1$) has budget~$2$
and $Z\equiv 2\pmod{11}$. This reproduces
Corollary~\ref{cor:div_ext}(v) directly from the atlas.
\end{remark}

%% ====================================================================
\section{Empirical killer sets}\label{sec:killer_sets}
%% ====================================================================

The atlas classification (Section~\ref{sec:atlas}) together with the
Erd\H{o}s inequality at each $k\ge 1$ suggests an empirical
investigation: for each hard residue~$r$, is there a small set
$S_r\subseteq[1,K_{\max}]$ such that every candidate $N$ with
$N\equiv r\pmod M$ fails the Erd\H{o}s inequality at some $k\in S_r$?

\begin{definition}[Base-enriched offsets]\label{def:baseenriched}
For residue $r$, the base-enriched $k$-set is the set of $k$-values
whose affine form participates in the enriched tuple of
Theorem~\ref{thm:complete}; explicitly,
\[
  \begin{array}{l|l}
    r & \text{base-enriched } k\text{-set}\\ \hline
    0    & \{1,2,3,4,5,6,8,9,10,12,18,20,24\}\\
    1716 & \{1,2,3,4,5,6,8,9,10,12,15,18,20,24,30\}\\
    4862 & \{1,2,3,4,5,6,8,9,10,12,14,18,20,24\}\\
    16302 & \{1,2,3,4,5,6,8,9,10,12,18,20,24,30\}\\
    17160 & \{1,2,3,4,5,6,8,9,10,12,18,20,24,28,36\}
  \end{array}
\]
The \emph{killer universe} at residue~$r$ is $[1,100]$ minus the
base-enriched $k$-set.
\end{definition}

\begin{definition}[Strong near-miss]\label{def:strong_near_miss}
For residue~$r$ with enriched-tuple size~$m_r$, a sample $N=Mt+r$
is a \emph{strong near-miss} iff at least $m_r-2$ of the $m_r$
enriched-tuple forms are prime at~$N$. The corresponding threshold
values are $11,13,12,12,13$ for
$r\in\{0,1716,4862,16302,17160\}$ respectively.
\end{definition}

\begin{definition}[Killer set]\label{def:killer_set}
For residue~$r$, the \emph{minimum killer set} $S_r$ is the smallest
subset of the killer universe at $r$ such that for every observed
strong near-miss $N$ at~$r$, some $k\in S_r$ satisfies
$\tau(2520N-k)>k+2$.
\end{definition}

\begin{proposition}[Small-cover phenomenon at $r=0$]\label{prop:killer_sets}
Let $r=0$ with the strong threshold of
Definition~\ref{def:strong_near_miss} (primes $\ge 11$ of $13$
enriched forms). Freezing the empirical claims to the audited
$301$-record strong corpus used by
\texttt{data/two\_views\_analysis.json} and
\texttt{data/robustness\_audits.json}, one has:
\begin{enumerate}
\item there exists a minimum killer set of size~$3$ covering
all $301$ records — specifically $\{15,16,27\}$ is the unique
3-element cover;
\item no $2$-element cover exists;
\item the monitor set
$S_{\mathrm{mon}}=\{7,11,13,16\}$ covers all $301/301$ records.
Across the combined strong-threshold corpus
($n=303$ pooled $r=0$, $r=1716$, $r=4862$), the monitor set
still covers every record, and the $\{15,16,27\}$ cover remains
universal on the $r=0$ slice.
\end{enumerate}
\end{proposition}

\begin{remark}[Exact minima are data-dependent; small-cover and monitor phenomena are stable]\label{rem:unstable_minima}
Three corpus-growth events demonstrate that specific minimum covers
shift while the small-cover phenomenon persists.  At an earlier
snapshot with $n=228$ strong near-misses, the unique size-$3$ cover
was $\{14,15,16\}$. At the $n=301$ snapshot, $\{14,15,16\}$ covered
$300/301$ (outlier at $t=3{,}336{,}277{,}668{,}150$), so the minimum
shifted to $\{15,16,27\}$.  We freeze the paper's empirical
statements at this audited $n=301$ / combined $n=303$ snapshot to
avoid mixing later telemetry into the theorem-facing text.  The
robust structural claim is not the
identity of any specific minimum cover but the persistence of
\emph{some} small cover together with $\{7,11,13,16\}$
monitor coverage on the frozen audited corpus.  We do not treat either $\{14,15,16\}$ or
$\{15,16,27\}$ as a theorem target; the theorem-shaped analog is a
branch-dichotomy (see Remark~\ref{rem:dichotomy}).
\end{remark}

\begin{remark}[Cross-residue data: pilot only]\label{rem:crossres_pilot}
At $r=1716$ and $r=4862$, the current strong-threshold corpus
contains $1$ record each (from Hetzner searches on $t\in[1,10^9]$).
At $r=16302$ and $r=17160$, the strong-threshold corpus is
currently empty. We do not state per-residue minimum killer sets
at these residues as formal propositions; the available data are
pilot-level only. The monitor-coverage claim in
Proposition~\ref{prop:killer_sets}(iii) includes these residues
and remains $1/1$ at each, consistent with (but not statistically
equivalent to) the $r=0$ result.
\end{remark}

\begin{remark}[Conditioning layer, diagnostic observation, and structural target]\label{rem:dichotomy}
The outlier $t=3{,}336{,}277{,}668{,}150$ is at the
\emph{strong near-miss} layer ($11$ of $13$ enriched forms prime;
enriched failures at $k=18,20$), not the \emph{true base-pass}
layer ($13$ of $13$ enriched forms prime). The distinct
conditioning layers form distinct stochastic regimes; the two
should not be conflated.

At the strong near-miss layer, branch classification of the $301$
records (\texttt{data/branch\_classification.json}) shows exactly
one record passes $\{14,15,16\}$ simultaneously; this
record is killed by $\{7,11,13,27\}$. Empirically at this layer,
\[
\mathrm{NearBasePass}_{11/13}(N)\wedge\bigwedge_{k\in\{14,15,16\}}\tau(2520N-k)\le k+2
\]
has observed density $1/301\approx 0.3\%$, and in the one observed
instance failure occurs at smaller offsets rather than at
$\{14,15,16\}$. This is a diagnostic observation at the strong
near-miss layer, not a theorem-shaped claim about true candidates.

At the true base-pass layer ($13/13$), the observed hit at
$t=11{,}741{,}230{,}040{,}310$ is killed by $\{14,15,16\}$ at every
offset. The outlier gives no direct evidence at this layer.

The natural structural target is therefore not a finite-cover
statement but a \emph{conditioned divisor-distribution model}: for
each omitted offset $k$ and each prime $\ell$, the local
probability
\[
\rho_{\ell,k}\;:=\;\Pr\bigl(\ell\mid Z_{0,k}(t)\,\big|\,\mathcal{B}\text{ locally passes mod }\ell\bigr),
\]
where $\mathcal{B}$ denotes the enriched $r=0$ base system and
$Z_{0,k}$ is the primitive cofactor from the atlas of
Section~\ref{sec:atlas}, together with the pairwise correlations
$\mathrm{Corr}(\omega(Z_{0,k}),\omega(Z_{0,j}))$. By an Euler-product
over $\ell\le Y$, this determines a conditioned
Erd\H{o}s--Kac-type distribution for $\omega(Z_{0,k}(t))$, whose
marginal and joint tails govern the Erd\H{o}s-inequality pass
probabilities. We leave the computation and consequences of this
model to future work; the outlier is one calibration point.

We do not promote $\{14,15,16\}$ as structurally sufficient. Its
observed strength ($99.67\%$ empirical dominance at the strong
near-miss layer) appears to come from high conditional failure
probability of omitted offsets with small divisor budgets, not
from forced incompatibility.
\end{remark}

\begin{remark}[Threshold normalization]
The strong-threshold choice (miss-count $\le 2$ relative to the
enriched tuple) is important: at looser thresholds the minimum
killer set becomes polluted by offsets $k$ that correspond to
base-enriched forms, artifacts of the incomplete base constraint
rather than genuine killers. In particular, at threshold
$\ge 10/15$ at $r=1716$, the minimum cover $\{1,13,42,72\}$
includes $k=1$; since $2520N-1$ is already one of the $15$
enriched forms at $r=1716$, samples failing at $k=1$ are precisely
samples that are not candidates at all. Filtering the killer
universe to exclude base-enriched $k$-values and restricting to
strong near-misses yields the clean estimates of
Proposition~\ref{prop:killer_sets}.
\end{remark}

\begin{remark}[Status: empirical, not theorem]\label{rem:emp_only}
Proposition~\ref{prop:killer_sets} is an empirical statement about
observed records. We tested three finite-incompatibility programs
for a proof of ``every candidate at $r=0$ fails the Erd\H{o}s
inequality at some $k\in\{14,15,16\}$'': a local-covering solver
at primes $p\le 200$ not in $P_0$; an extension adding core-form
primality constraints at $p\le 1000$; a further extension adding
the pairwise-gcd identities
$\gcd(n-a,n-b)\mid |a-b|$ (for $a,b\in\{14,15,16\}$, these give
$\gcd=1$ or $\gcd\mid 2$). Across all $54$ atlas branch triples at
$r=0$, none of the three programs produced a contradiction. This
is a strong negative result on three specific finite-combinatorial
vectors; it does not rule out other finite proofs, and it locates
the natural next step for a proof in the analytic direction — a
density-level control on the budget-$\ge 4$ Dickson-set tail,
with a tuple-restricted Titchmarsh-type bound as the nearest
proxy target, which we discuss in
Section~\ref{sec:conditional}. See
\texttt{docs/theorem\_sprint\_conclusion.md} for the full record.
\end{remark}

\begin{remark}[Triage pipeline and trigger events]\label{rem:triage}
Proposition~\ref{prop:killer_sets} is maintained as a live empirical
claim by a four-stage triage pipeline applied to every new
witness-search record: (i) base-enriched filter (drop any record
failing the enriched tuple at its residue, i.e., below strong
threshold); (ii) monitor filter against $S_{\mathrm{mon}}=\{7,11,13,16\}$;
(iii) full killer-universe scan over $[1,100]$ minus base-enriched
offsets; (iv) deep-survivor classification. A \emph{trigger event} is
any record that passes stages (i) and (ii) — a strong near-miss
that is not killed by the monitor. On the frozen audited combined
strong corpus (total $n=303$ records across $r=0$, $r=1716$, and
$r=4862$), the number of trigger events is zero: every strong-threshold
record at every residue fails the Erd\H{o}s inequality at some
$k\in S_{\mathrm{mon}}$. A single trigger event would require revisiting
the paper's empirical framing; continued null returns tighten it.
Details in \texttt{scripts/analyze\_two\_views.py} and
\texttt{docs/holding\_pattern\_rules.md}.
\end{remark}

\subsection{Branchwise effective budgets explain the small-offset obstruction profile}\label{subsec:beff_mechanism}

We now identify the structural mechanism governing the small-offset
obstruction profile documented above. The claim is that the empirical
pass/fail behavior at each offset $k$ is controlled by an
atlas-derived quantity, the \emph{branchwise effective budget}, and
that this quantity predicts the pass-rate ordering across the entire
killer universe $k\le 100$ at $r=0$ and explains the cross-residue
universality of the monitor set $\{7,11,13,16\}$.

\begin{definition}[Branchwise effective budget]\label{def:beff}
On each atlas branch $\lambda$ of $(r,k)$ (Section~\ref{sec:atlas}),
the atlas factorization
\[
  2520N - k \;=\; A_{r,k,\lambda}(t)\cdot Z_{r,k,\lambda}(t),
\]
with $A_{r,k,\lambda}(t) = \bigl(\prod_{p\notin S}p^{c_p}\bigr)\cdot\bigl(\prod_{p\in S}p^{a_p+\lambda_p}\bigr)$
the fixed--and--active divisor content and $Z_{r,k,\lambda}$ the
primitive cofactor coprime to $A = 2520M$, defines the
\emph{branchwise effective budget}
\[
  B^{\mathrm{eff}}_{r,k,\lambda}
  \;:=\;
  \bigl\lfloor (k+2)/\tau(A_{r,k,\lambda})\bigr\rfloor.
\]
\end{definition}

\begin{lemma}[Branch criterion for Erd\H{o}s pass/fail]\label{lem:branch_criterion}
For an Erd\H{o}s $\#647$ candidate $N = Mt + r$ lying on atlas branch
$\lambda$ at $(r,k)$, the Erd\H{o}s inequality
$\tau(2520N - k)\le k+2$ holds if and only if
\[
  \tau\bigl(Z_{r,k,\lambda}(t)\bigr) \;\le\; B^{\mathrm{eff}}_{r,k,\lambda}.
\]
In particular: $B^{\mathrm{eff}}=0$ forces automatic failure;
$B^{\mathrm{eff}}=1$ forces $Z=1$, impossible for sufficiently
large $N$; $B^{\mathrm{eff}}=2$ forces $Z$ prime or $Z=1$;
$B^{\mathrm{eff}}=3$ forces $Z$ prime, $Z = p^2$, or $Z=1$;
$B^{\mathrm{eff}}\ge 4$ admits multiple squarefree and prime-power
shapes.
\end{lemma}

\begin{proof}
Immediate from multiplicativity of $\tau$ applied to the atlas
factorization in Definition~\ref{def:beff} together with
$\gcd(A_{r,k,\lambda}, Z_{r,k,\lambda}) = 1$.
\end{proof}

The first-order conditioned Euler-product model
(computational check in
\texttt{scripts/conditioned\_divisor\_model.py}) shows that the
primitive cofactors $Z_{r,k}$ are nearly exchangeable under local
base conditioning at ordinary primes: one has
$\mu_k(Y=10^4)\approx 1.23$ and $\sigma_k\approx 1.10$ uniformly for
the test offsets $k\in\{7,11,13,14,15,16,27\}$ at $r=0$. Thus the
observed asymmetry of the small-offset obstruction profile is not
explained by first-order local divisibility. The asymmetry enters
through the branchwise fixed content $A_{r,k,\lambda}$ in
Definition~\ref{def:beff}.

\begin{proposition}[Audited effective-budget distribution at $r=0$]\label{prop:beff_r0}
For the $301$-record strong-near-miss corpus at $r=0$ (primes
$\ge 11$ of the $13$ enriched forms), the branchwise effective-budget
distribution at each test offset is:

\begin{center}
\renewcommand{\arraystretch}{1.05}
\begin{tabular}{c|rrrrrrr|c}
\hline
$k$ & $B^{\mathrm{eff}}{=}1$ & $2$ & $3$ & $4$ & $5$ & $6$ & $7$ & total \\
\hline
$7$ & $3$ & $7$ & $32$ & $259$ & $0$ & $0$ & $0$ & $301$ \\
$11$ & $0$ & $0$ & $0$ & $23$ & $0$ & $278$ & $0$ & $301$ \\
$13$ & $0$ & $0$ & $0$ & $0$ & $16$ & $0$ & $285$ & $301$ \\
$14$ & $1$ & $53$ & $0$ & $247$ & $0$ & $0$ & $0$ & $301$ \\
$15$ & $5$ & $161$ & $0$ & $135$ & $0$ & $0$ & $0$ & $301$ \\
$16$ & $7$ & $59$ & $234$ & $0$ & $0$ & $0$ & $0$ & $300^\ast$ \\
$27$ & $0$ & $0$ & $5$ & $27$ & $71$ & $0$ & $198$ & $301$ \\
\hline
\end{tabular}
\end{center}
\noindent{\footnotesize $^\ast$ one additional record at $k=16$ has
$B^{\mathrm{eff}}=0$ — fixed-content divisor count already exceeds
$k+2$, so Erd\H{o}s inequality automatically fails on that branch.}

In this corpus: every branch at $k=16$ has $B^{\mathrm{eff}}\le 3$
(so $Z$ must lie in $\{1, \text{prime}, p^2\}$);
$55\%$ of branches at $k=15$ have $B^{\mathrm{eff}}\le 2$
(so $Z$ must lie in $\{1, \text{prime}\}$);
$18\%$ of branches at $k=14$ have $B^{\mathrm{eff}}\le 2$;
and no branches at $k\in\{11,13\}$ have $B^{\mathrm{eff}}\le 4$.
\end{proposition}

\begin{proposition}[Severity ordering across $k\le 100$ at $r=0$]\label{prop:beff_ordering}
Let $p_k^{\mathrm{pred}}$ denote the branch-averaged prediction
\[
  p_k^{\mathrm{pred}}
  \;=\;
  \sum_{\lambda} w_{k,\lambda}\cdot
  F\bigl(B^{\mathrm{eff}}_{r,k,\lambda}\bigr),
\]
where $F$ is a Landau-order $\tau$-pass approximation and $w_{k,\lambda}$
is a branch-weight measure. Let $p_k^{\mathrm{emp}}$ denote the
empirical pass rate on the $301$-record strong corpus at $r=0$.
Across all $87$ offsets $k\in[1,100]$ not in the $r=0$
base-enriched set, the Spearman rank correlations are:
\[
  \rho(p^{\mathrm{pred}}, p^{\mathrm{emp}})
  \;=\;
  0.9386\ \text{(empirical branch weights)},
  \quad
  0.9330\ \text{(neutral-sampled branch weights)}.
\]
Empirical and neutral predictions themselves agree at
$\rho = 0.9780$. Chronological splitting of the corpus into halves
gives between-halves $\rho = 0.9639$, with neutral-prediction-vs-half
correlations of $0.9385$ and $0.9205$ respectively.
A bootstrap of $200$ resamples for the neutral-prediction audit gives
median $\rho = 0.9301$ with
$95\%$ confidence interval $[0.9218,\ 0.9375]$.
\end{proposition}

\begin{remark}[Scope of validation]
Proposition~\ref{prop:beff_ordering} validates the effective-budget
mechanism at the $\mathcal{N}_{11/13}$ strong-near-miss layer across
all offsets $k\le 100$ in the frozen audited $r=0$ corpus, under both
empirical and neutral branch weights, and robustly to chronological
partitioning and bootstrap. At the true base-pass layer
$\mathcal{B}_{13/13}$ only the single confirmed witness at
$t=11{,}741{,}230{,}040{,}310$ exists, so this layer is not
independently validated. The absolute magnitudes of the Landau
prediction are approximately a factor of $1.5$--$2$ below empirical
rates, reflecting the crudeness of the first-order tail law; the
ordering is stable. The main residuals occur at $k=30,36$, where a
$B^{\mathrm{eff}}=4$ branch with surviving pattern $[1,1]$ (semiprime)
dominates the empirical distribution while the Landau approximation
underweights the semiprime contribution. This indicates that
effective budget is the first-order mechanism; pattern-level
cofactor modeling is the natural second-order correction.
\end{remark}

\begin{proposition}[Cross-residue atlas tightness of $k=16$]\label{prop:k16_universal}
Neutral-sampled effective-budget statistics across all five hard
residues give the following values for $k=16$:
\begin{center}
\renewcommand{\arraystretch}{1.05}
\begin{tabular}{c|cccc}
\hline
$r$ & mean $B^{\mathrm{eff}}$ & fraction with $B^{\mathrm{eff}}\le 3$ & rank among $87$-$86$ non-base offsets \\
\hline
$0$      & $3.38$ & $49\%$ & $1$ \\
$1716$   & $3.35$ & $51\%$ & $3$ \\
$4862$   & $3.35$ & $51\%$ & $2$ \\
$16302$  & $3.37$ & $49\%$ & $2$ \\
$17160$  & $3.37$ & $49\%$ & $3$ \\
\hline
\end{tabular}
\end{center}
At every hard residue, $k=16$ lies among the three atlas-tightest
offsets. The $k=16$ atlas branches at each residue have no
$B^{\mathrm{eff}}>4$ entries (the active prime is $2$ with $a_2 = 3$,
constraining the branch exponent contribution to a narrow range).
The independence-product pass-rate scan of
Proposition~\ref{prop:case_a_density} independently identifies
$k=16$ as a top-3 bottleneck at every residue: empirical per-$k$
pass rates at $k=16$ are $0.336, 0.358, 0.351, 0.351, 0.348$
for $r\in\{0,1716,4862,16302,17160\}$ respectively, confirming the
atlas-geometric prediction at the observational level.
\end{proposition}

\begin{remark}[Moving-$K$ window empirical scan]\label{rem:moving_k_empirical}
Per-$k$ empirical pass rates at $n=1{,}000$ samples per $(r, k)$
(Hetzner cpx51, 16-worker parallel run) reveal that the tight
obstruction is concentrated at $k\le 100$.  Mean per-$k$ pass rates
at each hard residue, by $K$-block:

\begin{center}
\begin{tabular}{c|cccc}
\hline
$r$ & $[1, 100]$ & $[101, 200]$ & $[201, 500]$ & $[501, 1000]$ \\
\hline
$0$     & $0.806$ & $0.958$ & $0.992$ & $0.999$ \\
$1716$  & $0.810$ & $0.954$ & $0.991$ & $0.998$ \\
$4862$  & $0.805$ & $0.959$ & $0.991$ & $0.999$ \\
$16302$ & $0.810$ & $0.956$ & $0.991$ & $0.999$ \\
$17160$ & $0.808$ & $0.954$ & $0.991$ & $0.998$ \\
\hline
\end{tabular}
\end{center}

\noindent Very-tight offsets (per-$k$ pass rate $<0.3$) appear only in
$[1, 100]$ ($0$--$2$ per residue) and \emph{never} in $[101, 1000]$.
The deepest rate in $[101, 1000]$ is at $k=120$ ($\approx 0.50$ at
every residue).  The fixed $k\le 100$ screen captures the
\emph{sharp marginal} obstruction pressure; the moving tail at
$k>100$ does not introduce new individual lethal offsets.

\emph{Joint block survival.} A complementary per-$t$ experiment
(Hetzner, $n=500$ samples per residue, $16$-worker parallel)
tests the joint event $\bigcap_k \{\text{pass at }k\}$ within each
$K$-block:

\begin{center}
\begin{tabular}{c|ccccc}
\hline
$K$-block & $r=0$ & $r=1716$ & $r=4862$ & $r=16302$ & $r=17160$ \\
\hline
$[1, 100]$     & $0.000$ & $0.000$ & $0.000$ & $0.000$ & $0.000$ \\
$[101, 200]$   & $0.006$ & $0.004$ & $0.008$ & $0.004$ & $0.012$ \\
$[201, 500]$   & $0.070$ & $0.040$ & $0.040$ & $0.040$ & $0.054$ \\
$[501, 1000]$  & $0.526$ & $0.422$ & $0.456$ & $0.470$ & $0.348$ \\
\hline
\end{tabular}
\end{center}

\noindent Marginal per-$k$ rates near $1$ compound to meaningful
joint suppression: mean per-$k$ rate $0.958$ at $[101, 200]$ gives
naive product $0.958^{100}\approx 0.014$, matching observed
$0.004$--$0.012$; rate $0.999$ at $[501, 1000]$ gives $0.999^{500}\approx 0.61$,
matching observed $0.35$--$0.53$.  \emph{Cumulative}
$P(\text{pass all }k\in[1, K])$ is $0/500$ at every $K\in\{100,
200, 500, 1000\}$ across all residues.

\emph{Interpretation.}  The case-A screen at $k\le 100$ captures
the sharp marginal obstruction pressure; the moving tail at
$k>100$ contributes compound multiplicative suppression via
many mild constraints, not via new individual lethal offsets.
Both components are empirically non-trivial.  See
\texttt{scripts/moving\_k\_window\_experiment.py},
\texttt{scripts/joint\_block\_survival\_experiment.py},
\texttt{data/joint\_block\_survival\_experiment.json}, and
\texttt{docs/joint\_block\_survival\_result.md}.
\end{remark}

\begin{remark}[$k=14$ is a second stratum-invariant universal wall]\label{rem:k14_wall}
An analogous stratum-invariance analysis to Remark~\ref{rem:k7_stable}
reveals that $k=14$ shares the structural-wall property of $k=7$:

\begin{center}
\begin{tabular}{c|ccc}
\hline
stratum & sampled $n$ & $k=14$ pass rate & dominant $B^{\mathrm{eff}}$ \\
\hline
$C_{\ge 9}$  & $500$ & $31.6\%$ & $B=4$ ($85\%$) \\
$C_{\ge 10}$ & $500$ & $39.4\%$ & $B=4$ ($84\%$) \\
$C_{\ge 11}$ & $375$ & $33.9\%$ & $B=4$ ($83\%$) \\
$C_{\ge 12}$ & $23$  & $21.7\%^*$ & $B=4$ ($70\%$) \\
$C_{\ge 13}$ & $1$   & $0\%^*$  & $B=2$ \\
Monitor escapes & $143$ & $32.9\%$ & $B=4$ ($87\%$) \\
\hline
\end{tabular}
\end{center}

\noindent Sieve-weighted ratio $R_{14}(y)$ for
$y\in\{7,11,17,29,47,97\}$ evaluates to
\[
  \{1.052, 1.052, 1.052, 0.994, 0.975, 0.983\},
\]
essentially flat at $1$.  Base-tuple quality does not raise $k=14$
survival.  When $k=14$ fails, dominant $Z_{14}$-shape is $3$-distinct-primes
($\approx 45\%$ of fails), same as $k=7$.

The joint structure is notable: at monitor-escape records, $k=14$ pass
rate is $32.9\%$, which equals the unconditional rate — confirming
$k=14$ is TRANSVERSE to monitor passage.  Combined with the $k=7$
result, we have TWO independent stratum-invariant universal walls.  See
\texttt{scripts/k14\_transversality.py}.
\end{remark}

\begin{remark}[Monitor-escape family structure and $k=14$ backup bottleneck]\label{rem:monitor_escape_family}
Of the $29{,}290$ records in the $r=0$ near-miss corpus
($C_{\ge 9}$), $143$ pass the monitor $\{7,11,13,16\}$.  Their
first-fail distribution concentrates heavily at $k=14$:

\begin{center}
\begin{tabular}{c|c}
\hline
first-fail $k$ & count (of $143$) \\
\hline
$14$ & $96\ (67.1\%)$ \\
$15$ & $38\ (26.6\%)$ \\
$21$ & $4$ \\
$17$ & $3$ \\
$19, 22$ & $1$ each \\
\hline
\end{tabular}
\end{center}

\noindent At first-fail, the $B^{\mathrm{eff}}$ distribution is
dominated by $B=4$ ($97/143 = 67.8\%$) and $B=2$ ($36/143 = 25.2\%$).
At $B=2$ the branch forces $Z$ prime and the record fails because
$Z$ is semiprime or larger; at $B=4$ the record's $Z$ is a
$3$- or $4$-prime integer exceeding budget $\tau\le 4$.  The
dominant $Z$-shape at first-fail is $3$-distinct-primes
($65/143 = 45.5\%$).

\emph{Secondary monitor coverage.}  Across the $143$ monitor-escape
records, empirical catch rates are:
\begin{itemize}
\item $\{14, 15, 27\}$: $138/143 = 96.5\%$.
\item $\{15, 16, 27\}$: $129/143 = 90.2\%$.
\item $\{14, 16\}$: $96/143 = 67.1\%$.
\item $\{27\}$ alone: $87/143 = 60.8\%$.
\end{itemize}

\noindent The combined network $\{7,11,13,14,15,16,27\}$ catches
$138/143$ monitor-escapes at low $k$; the remaining $5$ pass the
entire network and first-fail at slightly higher $k$ (primarily
$k=17$, $k=21$, or at $k\in[101, 1000]$).  The emergent structural
statement is therefore:
\emph{the robust monitor is a low-offset network, not any single
minimum cover}.  Specific $3$-element covers continue to drift with
corpus growth; the network $\{7,11,13,14,15,16,27\}$ has been stable
across the $n=301,330,375$ snapshots.  See
\texttt{scripts/monitor\_escape\_family\_classification.py},
\texttt{scripts/batch\_autopsy\_143\_escapes.py} and
\texttt{data/batch\_autopsy\_143\_escapes.json}.
\end{remark}

\begin{remark}[Moving-tail is MORE suppressive conditional on monitor-escape]\label{rem:conditional_tail}
The unconditional moving-tail block survival of
Remark~\ref{rem:moving_k_empirical} measures
$\Pr(\text{pass block}\mid\text{raw } t)$.  The sharper conditional
question is:
\[
  \Pr(\text{pass block }[K_{\text{lo}}, K_{\text{hi}}]\mid \text{passes monitor }\{7,11,13,16\}).
\]
Using the $143$ observed monitor-escape records as a conditional
sample:

\begin{center}
\begin{tabular}{c|ccc}
\hline
$K$-block & raw $\Pr$ & $S_1$ (monitor-escape) $\Pr$ & ratio $S_1/{\rm raw}$ \\
\hline
$[101, 200]$  & $0.006$ & $0.000$ & $0.00$ \\
$[201, 500]$  & $0.070$ & $0.014$ & $0.20$ \\
$[501, 1000]$ & $0.526$ & $0.294$ & $0.56$ \\
\hline
\end{tabular}
\end{center}

\noindent Conditioning on monitor-escape makes the moving-tail block
survival \emph{lower}, not higher: by factor $0.00$--$0.56$ across
$K$-blocks.  Monitor-escape records are not tail-friendly; if anything,
they carry additional modular structure that suppresses survival at
$k>100$ as well.

Further, among the $143$ monitor-escape records, only $5$ pass the
full network
\[
  \{7,11,13,14,15,16,27\}.
\]
All $5$ first-fail at
$k\in\{17, 21\}$, all within $[1, 100]$; the full $k\le 100$ screen
therefore catches all $143$ monitor-escapes (no record in the
current corpus requires the moving tail for catching).  The
moving-tail analysis above provides additional multiplicative
suppression beyond the screen, not the first catch.  See
\texttt{scripts/conditional\_moving\_tail\_on\_escapes.py} and
\texttt{scripts/deep\_escape\_ledger.py}.
\end{remark}

\begin{remark}[Deep-network escape ledger and $k=21$ as a third wall]\label{rem:deep_escape_ledger}
The $5$ records passing the full network
$\{7,11,13,14,15,16,27\}$ are:

\begin{center}
\begin{tabular}{c|c|c|c|c}
\hline
$t$ & primes & first-fail $k$ & $\tau$ & bound \\
\hline
$256{,}333{,}196{,}208$ & $9$ & $21$ & $48$ & $23$ \\
$314{,}789{,}936{,}208$ & $9$ & $17$ & $192$ & $19$ \\
$508{,}200{,}324{,}660$ & $9$ & $17$ & $64$ & $19$ \\
$2{,}441{,}574{,}485{,}160$ & $9$ & $21$ & $32$ & $23$ \\
$56{,}750{,}755{,}674{,}480$ & $9$ & $21$ & $32$ & $23$ \\
\hline
\end{tabular}
\end{center}

\noindent The deep escapes concentrate at $k=21$ ($3/5$) and $k=17$
($2/5$).  Analogous transversality analysis to
Remarks~\ref{rem:k7_stable} and~\ref{rem:k14_wall} gives:

\begin{center}
\begin{tabular}{c|cccccc}
\hline
$k$ & $C_{\ge 9}$ & $C_{\ge 10}$ & $C_{\ge 11}$ & Monitor & Network & note \\
\hline
$7$  & $38.8\%$ & $36.8\%$ & $36.7\%$ & $36.5\%$ & $40\%$ & wall, $B=4$ \\
$14$ & $31.6\%$ & $39.4\%$ & $33.9\%$ & $32.9\%$ & $60\%$ & wall, $B=4$ \\
$21$ & $36.4\%$ & $36.6\%$ & $33.6\%$ & $30.8\%$ & $20\%$ & wall, $B=5$ \\
$17$ & $69.8\%$ & $68.8\%$ & $71.2\%$ & $71.3\%$ & $60\%$ & not a wall, $B=9$ \\
\hline
\end{tabular}
\end{center}

\noindent Three stratum-invariant walls: $k=7$, $k=14$, $k=21$, each
with pass rate $\approx 33$--$37\%$ and $B^{\mathrm{eff}}$ dominated
by $B=4$ or $B=5$.  $k=17$ is loose ($\sim 70\%$ pass rate,
$B^{\mathrm{eff}}=9$ dominant) and serves as an occasional
escape-catcher via rare $B=6$ branches rather than as a wall.

The extended low-offset network
$\{7,11,13,14,15,16,17,21,27\}$ catches all $143$ monitor-escapes and
all $5$ deep-network escapes in the current corpus.  See
\texttt{scripts/deep\_escape\_ledger.py} and
\texttt{scripts/k17\_k21\_transversality.py}.
\end{remark}

%% ====================================================================
%% L2A theorem — asymptotic wall thinning
%% ====================================================================

\begin{lemma}[Shape decomposition]\label{lem:shape_decomp}
Let $n\ge 1$.  If $\tau(n)\le 4$, then
\[
  n\in\{1,\,p,\,p^{2},\,p^{3},\,pq\},
\]
with $p,q$ prime and, in the last case, $p<q$.  Conversely, each of
these forms has $\tau(n)\le 4$.
\end{lemma}

\begin{proof}
Write $n=\prod_{i=1}^{r}p_{i}^{e_{i}}$.  Then
$\tau(n)=\prod(e_{i}+1)\le 4$ forces either $r=0$
($n=1$); or $r=1$ with $e_{1}\in\{1,2,3\}$ ($n\in\{p,p^{2},p^{3}\}$);
or $r=2$ with $e_{1}=e_{2}=1$ ($n=pq$, $p\ne q$).
\end{proof}

\begin{theorem}[L2A wall asymptotic thinning]\label{thm:l2a}
Fix an affine form $L(u)=\alpha u+\beta$ with $\alpha\ge 1$,
$(\alpha,\beta)=1$, and prime support of $\alpha$ contained in
$\{2,3,5,11,13,17,19\}$.  In particular the statement applies uniformly
to the three cases
$(\alpha,\beta)\in\{(8{,}314{,}020,\,-1),\,(16{,}628{,}040,\,-1),\,(5{,}542{,}680,\,-1)\}$,
which are the dominant-branch primitive cofactors at $k\in\{14,7,21\}$
respectively at the residue $r=0$, and to their residue-translated
analogues at the other four hard residues.  Let
\[
  S(U)\;:=\;\#\{u\in[1,U]:\tau(\alpha u+\beta)\le 4\}.
\]
Then
\[
  S(U)\;\ll_{\alpha,\beta}\;U\cdot\frac{\log\log U}{\log U}
  \qquad(U\to\infty),
\]
and in particular $S(U)=o(U)$.  There exists an effective
$U_{0}=U_{0}(\alpha,\beta)$ such that $S(U)\le U/2$ for every
$U\ge U_{0}$.  Explicitly, one may take $U_{0}=\lceil e^{300}\rceil
<2\cdot 10^{130}$ uniformly for the three listed $\alpha$.
\end{theorem}

\begin{proof}
Write $X:=\alpha U+|\beta|$.  By Lemma~\ref{lem:shape_decomp} every $u$
contributing to $S(U)$ satisfies $\alpha u+\beta\in\{1,p,p^{2},p^{3},pq\}$.
Accordingly
\[
  S(U)\;\le\;1+N_{p}(U)+N_{p^{2}}(U)+N_{p^{3}}(U)+N_{pq}(U),
\]
where $N_{\bullet}(U)$ counts $u\le U$ such that $\alpha u+\beta$
takes the corresponding shape.

\emph{Perfect-power branches.}  If $\alpha u+\beta=p^{2}$ or $p^{3}$,
then $\alpha u+\beta\le X$ forces $p\le X^{1/2}$ or $p\le X^{1/3}$
respectively, and each such $p$ determines at most one $u$.  Hence
\[
  N_{p^{2}}(U)\le X^{1/2}\ll_{\alpha,\beta}U^{1/2},
  \qquad
  N_{p^{3}}(U)\le X^{1/3}\ll_{\alpha,\beta}U^{1/3},
\]
both $o(U)$.

\emph{Prime branch.}  The condition $\alpha u+\beta=p$ is equivalent to
$p\in[2,\,\alpha U+\beta]$ and $p\equiv\beta\pmod{\alpha}$.  Since
$(\alpha,\beta)=1$ this is a reduced residue class.  By PNT in AP for
fixed modulus~$\alpha$,
\[
  N_{p}(U)=\frac{\operatorname{Li}(\alpha U)}{\varphi(\alpha)}
  +o_{\alpha}\!\left(\frac{U}{\log U}\right)
  \sim c_{\alpha}\frac{U}{\log U},
  \quad c_{\alpha}=\frac{\alpha}{\varphi(\alpha)}.
\]

\emph{Semiprime branch.}  If $\alpha u+\beta=pq$ with $p<q$, then
$p\le\sqrt{X}$, and $p\nmid\alpha$ (since $p\mid\alpha$ combined with
$p\mid pq=\alpha u+\beta$ would give $p\mid\beta$, contradicting
$(\alpha,\beta)=1$).  For each such~$p$ the remaining prime~$q$ lies
in the residue class $\beta p^{-1}\pmod{\alpha}$.  Hence
\[
  N_{pq}(U)\;\le\;\sum_{\substack{p\le\sqrt{X}\\p\nmid\alpha}}
  \pi\!\left(\frac{X}{p};\,\alpha,\,\beta p^{-1}\right).
\]
By PNT in AP uniformly over the finitely many reduced residue classes
$\bmod\,\alpha$,
\[
  \pi\!\left(\frac{X}{p};\alpha,a\right)
  =
  \frac{\operatorname{Li}(X/p)}{\varphi(\alpha)}
  +o_{\alpha}\!\left(\frac{X/p}{\log(X/p)}\right).
\]
For $p\le\sqrt{X}$ one has $\log(X/p)\ge\tfrac12\log X$, so
$\operatorname{Li}(X/p)\ll X/(p\log X)$.  Inserting and using Mertens
\[
  \sum_{\substack{p\le\sqrt{X}\\p\nmid\alpha}}\frac{1}{p}
  =\log\log X+B_{\alpha}+o(1)
\]
(with $B_{\alpha}=B_{1}-\sum_{p\mid\alpha}1/p$) yields
\[
  N_{pq}(U)\;\ll_{\alpha}\;
  \frac{X}{\varphi(\alpha)}\cdot\frac{\log\log X}{\log X}
  \;\asymp\;
  c_{\alpha}\,U\,\frac{\log\log U}{\log U}.
\]

\emph{Summation.}  Combining,
\[
  S(U)\;\le\;
  1+N_{p}(U)+N_{p^{2}}(U)+N_{p^{3}}(U)+N_{pq}(U)
  \;\ll_{\alpha,\beta}\;
  U\,\frac{\log\log U}{\log U}=o(U).
\]

\emph{Effective $U_{0}$.}  Replacing PNT in AP by the effective
Montgomery--Vaughan Brun--Titchmarsh inequality
$\pi(y;\alpha,a)\le 2y/(\varphi(\alpha)\log(y/\alpha))$, which holds
for every $y>\alpha$ with no exceptional-zero dependence, and
retracing Steps~1--3 at $\log U\ge 300$, one verifies that
$S(U)/U<0.463$ for all $U\ge e^{300}$ uniformly in the three listed
$\alpha$.  Thus $U_{0}=\lceil e^{300}\rceil<2\cdot 10^{130}$ suffices.
The constant is deliberately not optimised; sharper explicit
Brun--Titchmarsh or PNT-in-AP inputs (Bennett--Martin--O'Bryant--
Rechnitzer for $q\le 10^{5}$, Ramar\'e for wider ranges at
correspondingly larger $x_{0}$) lower $U_{0}$, but the resulting
threshold remains well beyond corpus scale.  This is the analytic
content of \emph{Level~2A}; the corpus-scale effective version
(\emph{Level~2B}) is discussed in Remark~\ref{rem:analytic_wall_program}.
\end{proof}

\begin{theorem}[L2B partial credit]\label{thm:l2b_partial}
Under the hypotheses of Theorem~\ref{thm:l2a}, the sharper envelope
\[
  S(U)\;\le\;c_{\alpha}\,U\,\frac{\log\log(\alpha U)+3}{\log(\alpha U)}
\]
holds for all $U\ge 2\cdot 10^{20}$.  Consequently
\[
  S(U)\;\le\;\tfrac{1}{2}U\qquad\text{for every }U\ge U_{0}^{\ast}
  \text{ with }
  U_{0}^{\ast}=6.6446042099\cdot 10^{24}.
\]
\end{theorem}

\begin{proof}[Proof sketch]
Replace classical Brun--Titchmarsh with the Yamada (2023) explicit
interval form
\[
  \pi(x+y;k,\ell)-\pi(x;k,\ell)
  <\frac{2y}{\varphi(k)(\log(y/k)+0.8601)},
  \qquad y>k,
\]
and retrace Steps~2--3 of the proof of Theorem~\ref{thm:l2a} with the
finite-range Mertens-in-AP bound
\[
  \left|\sum_{\substack{p\le y\\p\nmid\alpha}}\frac{1}{p}
    -\log\log y-B_{\alpha}\right|\le \varepsilon,\qquad
  \varepsilon=3.88\cdot 10^{-4}\;(y\ge 10^{5}),
\]
derivable by direct sieving up to $y=10^{8}$ and Dusart's
$|r(y)|\le 0.2/\log^{3}y$ for $y\ge 2{,}278{,}383$.  The prime
branch contributes
\[
  K_{\mathrm{prime}}(U)=\frac{2\log(\alpha U)}{\log U+0.8601},
\]
and the semiprime branch contributes
\[
  K_{\mathrm{semi}}(U)\le\log(\alpha U)\cdot T_{0.8601}(U)-\log\log(\alpha U),
\]
with $T_{\delta}(U)$ the Yamada-kernel hyperbola sum.  Direct
computation at $U=2\cdot 10^{20}$ gives
$K_{\mathrm{prime}}=2.6333$, $K_{\mathrm{semi}}\le 0.3434$,
$K_{\mathrm{cube}}<10^{-10}$, totalling $K\le 2.9767<3$.  Both
branches are monotonically decreasing for $U>2\cdot 10^{20}$.  The
threshold $U_{0}^{\ast}$ is the numerical root of
$c_{\alpha}(\log\log(\alpha U)+3)/\log(\alpha U)=1/2$.
Full numerical details are in
\texttt{docs/level2b\_partial\_credit\_codex.md}.
\end{proof}

\begin{remark}[Why the corpus-scale $U_{0}\le 10^{13}$ is still open]\label{rem:l2b_corpus_obstruction}
At $U=10^{13}$, the Yamada sharpening gives
$K_{\mathrm{prime}}=2.979$ and $K_{\mathrm{semi}}=1.149$, total
$K=4.128$, well above the threshold $K\le 0.75$ required for
$S(U)\le U/2$.  Maynard's $\eta$-gain for the Brun--Titchmarsh
constant is triggered only when $\log x/\log q\ge 8$, i.e.\
$x\ge \alpha^{8}$; for the prime branch this requires
$U\ge\alpha^{7}\approx 2.75\cdot 10^{48}$, and for the worst
semiprime inner range $U\ge\alpha^{15}\approx 10^{104}$.  Both
are out of reach at corpus scale.  Closing the corpus-scale
gap thus requires a technique beyond the Maynard/Yamada
Brun--Titchmarsh family --- see Remark~\ref{rem:analytic_wall_program}
for the surveyed alternatives.
\end{remark}

\begin{remark}[Analytic wall program: shape decomposition and rigour levels]\label{rem:analytic_wall_program}
The three stratum-invariant walls $W=\{7,14,21\}$ admit a clean
analytic reduction.  On the dominant branch $\lambda=(7:0)$, the
primitive cofactor satisfies
$Z_{r,k,(7:0)}(u)=\alpha u+\beta$ with
$\alpha=A/k$, $\gcd(\alpha,\beta)=1$.  The corresponding $B^{\mathrm{eff}}$
is $4$ for $k\in\{7,14\}$ and $5$ for $k=21$.  Passing the wall
requires $\tau(Z)\le B^{\mathrm{eff}}$.  Since Proposition~\ref{prop:b3_qnr_squarekill}
extends to the dominant $B^{\mathrm{eff}}=4$ branch by the same
residue obstruction on $\beta$ modulo the small primes, the shape
$Z=p^{2}$ is excluded, and the target reduces to the
\emph{shape-decomposition lemma}:
\[
  \tau(Z)\le 4,\;Z\ne p^{2}\;\Longrightarrow\;Z\in\{1,p,p^{3},pq\},
\]
which yields the explicit decomposition
\[
  S_{k}(U)\;\le\;\pi(\alpha U+|\beta|;\,\alpha,\beta)
  \;+\;(\alpha U)^{1/3}
  \;+\!\!\sum_{\substack{p\text{ prime}\\p\le\sqrt{\alpha U}\\\gcd(p,\alpha)=1}}\pi\!\left(\frac{\alpha U+|\beta|}{p};\,\alpha,\beta p^{-1}\right)\!.
\]

The remaining analytic question is the \emph{rigorous constant}
controlling the semiprime sum.  Numerical experiments
(\texttt{scripts/wall\_decomposition\_k14.py},
\texttt{scripts/wall\_empirical\_sample.py} and
\texttt{scripts/mertens\_correction\_test.py}) at
$\alpha=8{,}314{,}020$ give:
\begin{itemize}
  \item Prime-density heuristic matches empirical to $<3\%$; the
        Euler-product factor $c_{\alpha}=\alpha/\varphi(\alpha)\approx 5.01$
        is correct.
  \item The empirical semiprime constant at corpus scale fits at
        $c_{\mathrm{fit}}(U)\approx 3.32$, i.e.\
        $c_{\mathrm{fit}}/c_{\alpha}\approx 0.66$.  This \emph{is not}
        an asymptotic Selberg--Delange leading constant; it is the
        finite-range effect of excluding the small primes dividing
        $\alpha$ from the semiprime harmonic sum.  The cleanest
        statement uses the $\alpha$-restricted prime-harmonic
        function
        \[
          H_{\alpha}(y)\;:=\;\sum_{\substack{p\le y\\p\nmid\alpha}}\frac{1}{p}.
        \]
        Within a standard Mertens expansion
        $\sum_{p\le y}1/p=\log\log y+M_{\mathrm{Mertens}}+o(1)$, the
        subtraction $\sum_{p\mid\alpha}1/p=M_{\alpha}$ is dominant at
        our scales, giving the approximation
        \[
          \frac{c_{\mathrm{eff}}(U)}{c_{\alpha}}
          \;\approx\;
          1-\frac{M_{\alpha}}{\log\log(\alpha U)},
          \qquad
          M_{\alpha}\;=\;\sum_{p\mid\alpha}\frac{1}{p}
          \;=\;\tfrac12+\tfrac13+\tfrac15+\tfrac1{11}+\tfrac1{13}+\tfrac1{17}+\tfrac1{19}
          \;\approx\;1.313.
        \]
        The simple first-order prediction
        $1-M_{\alpha}/\log\log(\alpha U)$ sits roughly $6$--$9\%$
        \emph{below} the empirically-measured ratio across
        $U\in[10^{8},10^{17}]$ (measured on a full
        $k\times r\times\log_{10}U$ matrix at $10^{5}$ samples per
        cell using a Pollard-Rho shape classifier).  The directional
        agreement --- identical across the three walls
        $k\in\{7,14,21\}$ and matching the Mertens-correction
        prediction that all three walls share a structural
        constant --- is the important content.  The residual
        $6$--$9\%$ gap is expected: the simple formula captures the
        dominant $\alpha$-dependent subtraction $M_{\alpha}$ but
        omits lower-order Mertens/AP/boundary corrections.  The
        exact analytic object is the finite prime-harmonic sum
        $H_{\alpha}(\sqrt{\alpha U})$ inside a one-sided semiprime
        envelope; the first-order expression is a leading
        approximation.
  \item Classical Brun--Titchmarsh with $c_{\mathrm{mult}}=2$ gives
        $U_{0}\sim 10^{54}$ --- unusable at corpus scale.
  \item Under the empirical finite-range Mertens-corrected envelope,
        the decomposition upper bound drops below $U/2$ (both
        pointwise and cumulatively) at $U_{0}\approx 10^{8}$, well
        within the corpus scale.  The analytic target consists of an
        \emph{explicit one-sided upper envelope} for the semiprime
        sum of the form
        \[
          N_{pq}(U;\alpha,\beta)
          \;\le\;
          \frac{\alpha U}{\varphi(\alpha)}
          \sum_{\substack{p\le\sqrt{\alpha U}\\p\nmid \alpha}}
          \frac{1}{p\log(\alpha U/p)}
          \cdot(1+\text{explicit error})
        \]
        together with Brun--Titchmarsh-level bounds for $N_{p}$ and the
        trivial bound for $N_{p^{3}}$.  We do not claim this as a
        theorem; we state the target envelope and record that the
        empirical data support it at $U_{0}\approx 10^{8}$ for all
        three walls.
\end{itemize}

We structure the rigour roadmap into five sub-levels:
\begin{description}
  \item[\textbf{Level 1}] Empirical wall (achieved): pass rate
        $\approx 1/3$ at each $k\in W$, stratum-invariant to $\pm 4$~pp.
  \item[\textbf{Level 2A}] Unconditioned asymptotic thinning:
        $\#\{u\le U:\tau(\alpha u+\beta)\le 4\}=o(U)$.  Reachable by
        fixed-modulus prime-number-theorem-in-AP plus the shape
        decomposition.  Existence of $U_{0}$ follows but the explicit
        value is not controlled.
  \item[\textbf{Level 2B}] Effective corpus-scale thinning:
        $\#\{u\le U:\tau(\alpha u+\beta)\le 4\}\le U/2$ for all
        $U\ge U_{0}$ with explicit $U_{0}\le 10^{13}$.  Blocked by
        the Brun--Titchmarsh factor-$2$ in classical upper-bound
        sieves.  Using the character decomposition
        $\pi(y;\alpha,a)=\varphi(\alpha)^{-1}\sum_{\chi}\overline{\chi(a)}\pi_{\chi}(y)$,
        the principal contribution at $X=\alpha U\approx 8\cdot10^{19}$
        ($U=10^{13}$) is
        $P_{2}^{(0)}(X)\approx 3.13\cdot10^{12}$, already within the
        $c_{\mathrm{bt}}\le 1.25$ budget
        ($D_{\mathrm{semi}}\approx 2.75\cdot10^{12}$, giving
        principal-only $c_{\mathrm{eff}}^{(0)}=1.140$).  The theorem
        therefore reduces to bounding the nonprincipal contribution
        $|E_{\alpha}(X)|$ by $\approx 3.03\cdot10^{11}$.  We
        surveyed the standard candidate methods: large-sieve
        aggregate ($\approx 1.86\cdot10^{16}$; misses by
        $\sim\!6\cdot10^{4}$), P\'olya--Vinogradov / Burgess on
        primes ($\approx 3.83\cdot10^{13}$; misses by $\sim\!126$),
        dispersion method (no explicit finite-scale bound
        available), unconditional zero-density estimates (unable to
        rule out the exceptional-zero contribution at the required
        strength), and sharper Brun--Titchmarsh of Maynard and
        Yamada type (gives $\approx 8.8\cdot10^{12}$; misses by
        $\sim\!29$).  None supplies the required one-sided
        AP-semiprime envelope with constant $\lesssim 1.25$ at this
        $X$.  We therefore do not pursue Level~2B unconditionally
        in this paper.  A speculative bespoke approach --- an
        explicit weighted $\Lambda\!*\!\Lambda$ character expansion,
        or certified segmented sieving over the
        $\varphi(\alpha)=1{,}658{,}880$ invertible residue classes
        mod~$\alpha$, using the fact that
        $\alpha=2^{2}\cdot 3^{2}\cdot 5\cdot 11\cdot 13\cdot 17\cdot 19$
        has only $12$ primitive-conductor divisors exceeding
        Platt's $q\le 400{,}000$ verified range --- might yield the
        required finite constant as a computational project; we
        flag it as future work.
  \item[\textbf{Level 3}] Conditional wall lemma under the
        \emph{base-conditioned local equidistribution} hypothesis
        $\mathrm{BCLE}_{r}(Y)$: level-of-distribution $Y=U^{\theta}$
        for some $\theta>0$ on the base-tuple-conditioned set
        $\mathcal{B}_{r}(U)$ (no known plug-in theorem; cf.\ Maynard's
        dense-cluster hypothesis and Benatar's subset-BV framework;
        analogous in role to Elliott--Halberstam).
  \item[\textbf{Level 4}] Multi-wall cascade: combine Level-3 bounds
        at all three $k\in W$ with a controlled cross-offset
        correlation hypothesis; this is the level at which the
        escape-ladder heuristic approaches a conditional finiteness
        statement.
\end{description}

None of Levels~2--4 is claimed in this paper; the remark records
the target structure and the empirical constants that make it
plausible.  See \texttt{docs/analytic\_wall\_k14.md},
\texttt{docs/bcle\_hypothesis\_formalization.md} and
\texttt{docs/heuristic\_finiteness\_sketch\_v2.md} for the detailed
programme.
\end{remark}

\begin{remark}[$k=7$ is a stable bottleneck under conditioning]\label{rem:k7_stable}
A complementary empirical finding at the $r=0$ near-miss corpus:
the per-$k$ pass rate at $k=7$ is approximately constant across
conditioning strata $C_{\ge 9}, C_{\ge 10}, C_{\ge 11}$ (primes~$\ge 9,
10, 11$ of the $13$ enriched forms respectively).  Sample rates:

\begin{center}
\begin{tabular}{c|ccc}
\hline
stratum & sampled $n$ & $k=7$ pass rate & dominant $B^{\mathrm{eff}}$ at $k=7$ \\
\hline
$C_{\ge 9}$  & $500$ & $38.8\%$ & $B^{\mathrm{eff}}=4$ ($87\%$) \\
$C_{\ge 10}$ & $500$ & $36.8\%$ & $B^{\mathrm{eff}}=4$ ($83\%$) \\
$C_{\ge 11}$ & $330$ & $36.7\%$ & $B^{\mathrm{eff}}=4$ ($86\%$) \\
$C_{\ge 12}$ & $21$  & $57.1\%^*$ & $B^{\mathrm{eff}}=4$ ($90\%$) \\
$C_{\ge 13}$ & $1$   & $0\%^*$  & $B^{\mathrm{eff}}=4$ \\
\hline
\end{tabular}
\end{center}

\noindent $^*$Small-sample rates (binomial-compatible with the
uniform $\sim 37\%$ seen in large strata; not separately informative
at $n=21$ and $n=1$).  The $B^{\mathrm{eff}}$-branch distribution at
$k=7$ is also stable: dominated by $B^{\mathrm{eff}}=4$ with $\approx
10\%$ at $B^{\mathrm{eff}}=3$ and $\approx 2\%$ at $B^{\mathrm{eff}}=2$.
Mean $\tau(2520N-7)$ is $\approx 20$ across strata.  When $k=7$ fails,
the dominant $Z_7(t)$-shape is a $3$-distinct-prime integer ($\approx
50\%$ of fail cases), followed by $4$-prime and $5$-prime shapes.

This empirical stability of $k=7$ across conditioning explains why
the single confirmed $13$-tuple witness at
$t=11{,}741{,}230{,}040{,}310$ fails at $k=7$: the witness is
$C_{\ge 13}$, but the $k=7$ obstruction does not weaken with
stronger base-tuple primality.

A complementary sieve-weighted surrogate confirms the transversality:
define the base-quality score
$Q_y(t) = \#\{d\in\mathcal B : dN-1\text{ has no prime factor}\le y\}$,
where $\mathcal B$ is the $13$-form enriched tuple at $r=0$; sample
$n=2{,}000$ at $t\in[10^8,10^9]$ and compute the ratio
$R_7(y) = \Pr(k=7\text{ pass}\mid Q_y\ge Q_y^{(75\%)})/
\Pr(k=7\text{ pass}\mid\text{raw})$.
Staged across $y\in\{7,11,17,29,47,97\}$, $R_7(y)$ takes values
$\{0.971, 0.971, 0.971, 0.942, 0.936, 0.922\}$, all in $[0.92, 0.97]$.
High base-quality samples have, if anything, a slightly \emph{lower}
$k=7$ pass rate than raw; there is no evidence that base-tuple
conditioning raises $k=7$ survival.  This is a positive
falsification of the hypothesis that caveat (b) could close the
remaining gap in the density model at $k=7$.  See
\texttt{scripts/k7\_sieve\_weighted\_surrogate.py},
\texttt{data/k7\_sieve\_weighted\_surrogate.json},
\texttt{scripts/monitor\_escape\_k7\_autopsy.py},
\texttt{data/monitor\_escape\_k7\_autopsy.json}.
\end{remark}

\begin{remark}[$k=16$ at the strongest-near-miss subset]\label{rem:k16_strong_layer}
At the $r=0$ strong threshold with $12$ or more primes of the $13$
enriched forms ($n = 21$ records in the current corpus), $k=16$ alone
fails the Erd\H{o}s inequality in every case: $\tau(2520N-16) > 18$
for all $21$ records at this threshold.  At the broader $r=0$
strong-near-miss telemetry snapshot (primes $\ge 11$, $n = 375$), $k=16$ alone
kills $\approx 87\%$ of records, with the monitor $\{7, 11, 13, 16\}$
required to cover the remaining $44$.  The $21/21$ rate at the
$k\ge 12$ subset is consistent with — though a stronger observation
than — the atlas-tightness of
Proposition~\ref{prop:k16_universal}.  Sample size at the
$k\ge 12$ layer is small; we report this as an empirical observation
from a later telemetry snapshot, not as a universal claim.
\end{remark}

\begin{remark}[Structural explanation of the monitor $\{7,11,13,16\}$]\label{rem:monitor_structural}
The universality of the monitor $\{7,11,13,16\}$ decomposes into
residue-generic and residue-specific components:
\begin{enumerate}
\item \emph{Residue-generic:} $k=16$ is atlas-tight at every
residue (Proposition~\ref{prop:k16_universal}); $k=7$ ranks
within the top $8$--$10$ tightest offsets at every residue (mean
$B^{\mathrm{eff}}\approx 3.83$ across all $r$).
\item \emph{Residue-specific extreme-tightness:}
$k=13$ becomes extreme-tight at $r=1716$ and $r=16302$ (mean
$B^{\mathrm{eff}}\approx 2.85$, $100\%$ of branches with
$B^{\mathrm{eff}}\le 3$, rank $1$--$2$), reflecting the active-prime
$13$ content at those residues; $k=11$ becomes extreme-tight at
$r=17160$ (rank $2$, $100\%$ with $B^{\mathrm{eff}}\le 3$),
reflecting the analogous active-prime $11$ content.
\end{enumerate}
Thus at every hard residue, at least one member of $\{7,11,13,16\}$
lies in the top three atlas-tightest offsets. The monitor's
empirical coverage is consistent with this branch-budget geometry,
though the formal coverage claim remains empirical.
\end{remark}

\begin{remark}[Alternate universal-tight sets and empirical universality]\label{rem:alt_monitor}
Audit~4 identifies $\{7,16,40,72\}$ as the universal atlas-tight-top-$10$
intersection across all five hard residues — a candidate alternate
monitor family.  However, on the expanded strong-near-miss corpus
of $330$ records at $r=0$ (primes $\ge 11$, updated 2026-04-18),
$\{7,16,40,72\}$ misses one record at $t=2{,}445{,}307{,}576{,}020$
($99.70\%$ empirical coverage), while the empirically-motivated
monitor $\{7,11,13,16\}$ still achieves $100\%$.  A second
three-element universal monitor, $\{15,16,27\}$, also achieves
$100\%$ coverage on the expanded corpus and is a candidate
simplification; $\{7,11,13,16\}$ remains the base monitor used
throughout the paper because $k=16$ alone covers $86.67\%$ and
each of $\{7,11,13\}$ contributes residue-specific lethality
(see Remark~\ref{rem:monitor_structural}).  We flag these
alternatives for subsequent analysis; none replaces the
empirically-validated monitor here.
\end{remark}

\begin{remark}[Interpretation and scope]\label{rem:beff_interpretation}
The branchwise severity ordering of omitted offsets is a
theorem-shaped consequence of atlas arithmetic
(Lemma~\ref{lem:branch_criterion}); the empirical pass-rate
distribution at each offset reflects this ordering under a Landau
tail law for primitive cofactors. The empirical killer-set
phenomenon and the cross-residue monitor universality follow from
the branch-budget geometry. We do not claim that the monitor
$\{7,11,13,16\}$ provably covers every Erd\H{o}s candidate, nor
that the effective-budget mechanism constitutes a proof of
finiteness for $\#647$. What we claim is:
\begin{itemize}
\item the branch criterion of Lemma~\ref{lem:branch_criterion} is
exact;
\item the severity ordering it induces matches empirical pass-rate
ordering at $r=0$ across $87$ offsets with Spearman correlation
$0.93$--$0.94$ on the frozen audited corpus, robustly;
\item the monitor's coverage has a specific structural
explanation, though its formal universality remains an empirical
observation.
\end{itemize}
Computation scripts: \texttt{scripts/cofactor\_tail\_model.py},
\texttt{scripts/cofactor\_tail\_model\_v3.py},
\texttt{scripts/robustness\_audits.py},
\texttt{scripts/audit4\_cross\_residue.py}.
Data: \texttt{data/cofactor\_tail\_model.json},
\texttt{data/robustness\_audits.json},
\texttt{data/audit4\_cross\_residue.json}.
\end{remark}

\subsection{Rigorous-elimination layer and the survivor corridor}\label{subsec:beff_covering}

The effective-budget hierarchy of Lemma~\ref{lem:branch_criterion}
admits a natural deterministic/heuristic split at $B^{\mathrm{eff}}=1$.

\begin{proposition}[Automatic-failure layer is noncovering]\label{prop:ble1_noncover}
Let $\mathcal E_{r}^{\le 1}
:=\bigcup_{k\le 100,\;\lambda:\,B^{\mathrm{eff}}_{r,k,\lambda}\le 1}
  \mathcal C_{r,k,\lambda}$
denote the union of all atlas branch-constraint classes at which the
effective budget forces an automatic Erd\H{o}s failure.  Then for
every hard residue $r\in\{0,1716,4862,16302,17160\}$, the set
$\mathcal E_{r}^{\le 1}$ has positive complement in $\mathbb Z$.
\end{proposition}
\begin{proof}
For $r\in\{0,4862,16302,17160\}$, the measure-sum bound
\[
  \mathbb P(t\in\mathcal E_{r}^{\le 1})
  \le
  \sum_{(k,\lambda):\,B^{\mathrm{eff}}_{r,k,\lambda}\le 1}
  \frac{|\mathcal C_{r,k,\lambda}|}{\mathrm{cm}_{r,k,\lambda}}
\]
evaluates to $\{0.17,0.45,0.18,0.46\}$ respectively (computed over all
$(k,\lambda)$ with $B^{\mathrm{eff}}\le 1$), all strictly less than~$1$.
Pigeonhole then gives a positive-density complement.

For $r=1716$, the measure-sum bound is $1.14>1$, so the pigeonhole
argument does not directly apply.  We instead exhibit explicit
uncovered witnesses.  The $50$ values of $t$ in
\texttt{data/b\_le1\_uncovered\_certificates.json} at $r=1716$ (e.g.\
$t=314{,}985{,}525$, giving $N=14{,}548{,}866{,}415{,}941$;
$t=495{,}450{,}098$, giving $N=22{,}884{,}344{,}578{,}238$;
$t=422{,}888{,}998$, giving $N=19{,}532{,}819{,}930{,}338$) each have
minimum $B^{\mathrm{eff}}_{1716,k,\lambda(t)}\ge 2$ across all
$k\in[1,100]\setminus K_{\mathrm{base-enriched}}(1716)$, verified
independently by \texttt{scripts/verify\_r1716\_uncovered\_witness.py}.
Each such witness certifies that its arithmetic progression
$\{N\equiv 1716\pmod M: t\ge t_0\}$ lies outside $\mathcal E_{1716}^{\le 1}$.
\end{proof}
\begin{remark}[Quantitative escape rates]\label{rem:ble1_mc}
Monte Carlo estimates of $\mathbb P(t\in\mathcal E_{r}^{\le 1})$ at
$N=10^5$ random $t\in[0,L_r)$, where $L_r$ is the lcm of branch moduli,
give per-residue coverage fractions in the range $10$--$38\%$ (see
\texttt{scripts/covering\_analysis.py}).  The rigorous automatic-failure
layer does not by itself imply finiteness of Erd\H{o}s candidates at any
hard residue; for that, the sparse-shape layer $B\in\{2,3\}$ is
required (Proposition~\ref{prop:b3_qnr_squarekill} and
Remark~\ref{rem:survivor_corridor}).
\end{remark}

\begin{remark}[Survivor corridor and the $B\le 3$ sparse-shape layer]\label{rem:survivor_corridor}
The Erd\H{o}s candidate set at a hard residue $r$ survives
$\mathcal E_{r}^{\le 1}$ with probability $10\%$--$90\%$
(Proposition~\ref{prop:ble1_noncover}); the survivors are further
constrained by the next tightness layer
$B^{\mathrm{eff}}\in\{2,3\}$, each of which forces the primitive
cofactor $Z_{r,k,\lambda}(t)$ to lie in $\{1,p,p^2\}$ (or, after
Proposition~\ref{prop:b3_qnr_squarekill}, in $\{1,p\}$ at $k\le 100$).
Conditioning further on actual-primality of these shape-forced $Z$'s,
Monte Carlo at $N=3{,}000$ random $t\in[10^8,10^9]$ per residue
gives ``E3Q survivor'' rates of $2.00\%, 0.17\%, 0.33\%, 0.53\%, 0.10\%$
for $r\in\{0,1716,4862,16302,17160\}$, reductions of $50$--$1600\times$
from the $B^{\mathrm{eff}}\le 1$ rates.  The distribution of forced-prime
constraint counts among E3Q survivors per residue is:

\begin{center}
\begin{tabular}{c|rrrrrr|r}
\hline
$r$ & $0$ & $1$ & $2$ & $3$ & $4$ & $5$ & total \\
\hline
$0$      & $19$ & $19$ & $19$ & $2$ & $1$ & $0$ & $60$ \\
$1716$   & $0$  & $0$  & $0$  & $4$ & $0$ & $1$ & $5$ \\
$4862$   & $0$  & $3$  & $5$  & $2$ & $0$ & $0$ & $10$ \\
$16302$  & $0$  & $6$  & $5$  & $2$ & $2$ & $1$ & $16$ \\
$17160$  & $0$  & $0$  & $0$  & $2$ & $1$ & $0$ & $3$ \\
\hline
\end{tabular}
\end{center}

\noindent At $r=0$, $31.7\%$ of E3Q survivors carry \emph{zero} forced-prime
$B\le 3$ constraints: they land on $B^{\mathrm{eff}}\ge 4$ branches at
every omitted $k$ and bypass the forced-prime layer entirely.  These
``loose corridors'' dominate the residual E3Q density at $r=0$ and
$r\in\{4862,16302\}$, and must be modeled separately from the tight
corridors at $r\in\{1716,17160\}$ where every E3Q survivor carries
$\ge 3$ forced-prime constraints.  See
\texttt{scripts/e3q\_forced\_prime\_histogram.py} and
\texttt{data/e3q\_forced\_prime\_histogram.json}.
\end{remark}

\begin{remark}[Structural primitivity of the cofactor form]\label{rem:Z_primitive}
At every atlas branch $(r,k,\lambda)$, the primitive cofactor
$Z_{r,k,\lambda}(u) = \alpha\cdot u + \beta$ has
$\gcd(\alpha,\beta)=1$ where the slope
\[
  \alpha = \prod_{p\not\in S_{r,k}} p^{a_p - c_p(r,k)}
\]
is a product of the non-active small primes (with $S_{r,k}$ the
active-prime set and $a_p=v_p(A_{\mathrm{const}})$, $c_p(r,k)=v_p(C_{r,k})$,
and $a_p > c_p(r,k)$ for $p\not\in S_{r,k}$).  Empirically verified at
all $1{,}175$ branches with $B^{\mathrm{eff}}=2$ and all $760$ branches
with $B^{\mathrm{eff}}=3$ at $k\le 100$, and agreeing with the
structural derivation.  Consequently $Z$ admits no fixed-divisor
obstruction beyond the atlas factorization itself, ruling out the
simplest class of rigorous branch contradictions.  See
\texttt{scripts/b23\_shape\_contradiction\_scan.py} and
\texttt{data/b23\_shape\_contradictions.json}.
\end{remark}

\begin{proposition}[QNR square-kill certificate theorem at $k\le 100$]\label{prop:b3_qnr_squarekill}
For every atlas branch $(r,k,\lambda)$ with $r\in\{0,1716,4862,16302,17160\}$,
$k\in[1,100]\setminus K_{\mathrm{base-enriched}}(r)$, and
$B^{\mathrm{eff}}_{r,k,\lambda}=3$, there exists a prime $q$ dividing
$\alpha_{r,k,\lambda}$ such that $\beta_{r,k,\lambda}$ is a quadratic
non-residue modulo $q$ (or $q=2$ with $v_2(\alpha)\ge 2$ and
$\beta\not\equiv 1\pmod 4$).  Consequently $Z_{r,k,\lambda}(u)$ is
never a perfect square, and the branch shape requirement
$\tau(Z)\le 3$ collapses to ``$Z$ prime or $Z=1$''.
\end{proposition}

\begin{proof}
Finite certificate verification.  For each of the $1{,}242$ branches
$(r,k,\lambda)$ with $B^{\mathrm{eff}}_{r,k,\lambda}=3$ at
$k\le 100$, the accompanying data file
\texttt{data/b3\_qnr\_certificates\_k100.json} records
$(\alpha,\beta,q,\beta\bmod q,(\beta/q))$ with either
(i) $q$ an odd prime, $q\mid\alpha$, $(\beta/q)=-1$, or
(ii) $q=2$, $v_2(\alpha)\ge 2$, $\beta\bmod 4\ne 1$ (or, when
$v_2(\alpha)\ge 3$, $\beta\bmod 8\ne 1$).  In case (i),
$Z(u)\equiv\beta\pmod q$ is a quadratic non-residue for all $u$, so
$Z$ is not a perfect square.  In case (ii), $p^2\equiv 1\pmod 4$
(resp.\ $\pmod 8$) for all odd primes $p$, so $Z\ne p^2$.  Since
$\gcd(\alpha,\beta)=1$ by Remark~\ref{rem:Z_primitive}, $q\nmid\beta$,
closing the $p=q$ loophole.  Exhaustiveness (saturation lemma).  If $B^{\mathrm{eff}}_{r,k,\lambda}=3$
then $(k+2)/\tau(A_{r,k,\lambda})\ge 3$, so
$\tau(A_{r,k,\lambda})\le (k+2)/3\le 34$ for $k\le 100$.
Each active prime $p\in S$ contributes a factor $a_p+\lambda_p+1\ge 2$
to $\tau(A_{r,k,\lambda})$, so
$a_p+\lambda_p+1\le 34$ for every $p\in S$, forcing $\lambda_p\le 32$.
The enumeration of $B=3$ branches at cutoff $\lambda\le 50$ is therefore
rigorously saturated, and indeed the branch count agrees at cutoffs
$\lambda\le 30$ and $\lambda\le 50$ (both give $1{,}242$). See
\texttt{scripts/b3\_qnr\_certificate\_k100.py} and
\texttt{scripts/verify\_b3\_qnr\_certificates\_k100.py}.
\end{proof}

\begin{remark}[Extension at $k\le 1000$]\label{rem:b3_qnr_k1000}
The analogous scan at $k\le 1000$ produces QNR-certificate coverage
$\{100.00, 98.20, 98.83, 98.79, 97.77\}\%$ for
$r\in\{0,1716,4862,16302,17160\}$; overall $98.70\%$ of $12{,}495$
branches.  The $1.30\%$ residual is consistent with the
$2^{-\omega(\alpha)}$ random-QR exception rate given the typical
number of distinct prime factors of $\alpha$.  The extension is
therefore an empirical density regularity, not a finite-certificate
theorem.  Scripts and data:
\texttt{scripts/b23\_qnr\_deep.py}, \texttt{data/b23\_qnr\_deep.json}.
\end{remark}

\begin{remark}[Interpretation: where a deterministic proof can still live]\label{rem:beff_frontier}
Remark~\ref{rem:survivor_corridor},
Proposition~\ref{prop:b3_qnr_squarekill} and
Remark~\ref{rem:b3_qnr_k1000} together describe the current proof
frontier.  The $B^{\mathrm{eff}}\le 1$ layer is deterministic but
non-covering.  The $B^{\mathrm{eff}}=2,3$ layer is uniformly
shape-collapsed to ``$Z$ prime or $Z=1$'' at $k\le 100$ (by
Proposition~\ref{prop:b3_qnr_squarekill}), and heuristically
collapsed at $98.70\%$ of $B^{\mathrm{eff}}=3$ branches at $k\le 1000$;
but the collapsed layer does not yield a rigorous per-branch
contradiction (since ``$Z$ prime'' is not structurally excluded).  A deterministic-finiteness route from this framework
would need to combine (i) pairwise $\gcd$-at-$t$ constraints on
simultaneous prime cofactors, (ii) covering-system arguments on the
linear forms $Z_{r,k,\lambda}$, or (iii) shape obstructions beyond
QNR (cubic/quartic characters).  The current paper does not pursue
these; conditional thinning statements follow from
Proposition~\ref{prop:beff_ordering} combined with Bateman--Horn on
the resulting shape conditions.
\end{remark}

%% ====================================================================
\section{Universal admissibility}\label{sec:admissibility}
%% ====================================================================

\begin{theorem}[Universal admissibility of the core sub-tuple, Lean-verified]
For each $r\in\{0,1716,4862,16302,17160\}$, the $13$-form core
sub-tuple \texttt{hardRootCoeffs} is admissible at every prime~$p$.
\end{theorem}

\begin{proof}
For $p>13$: each of the $13$ forms contributes at most one forbidden
residue mod~$p$; pigeonhole gives a surviving residue.
For $p\le 13$: discharged by \texttt{native\_decide} independently
at each hard residue (each generating its own computational
verification axiom; no sharing between residues).
\end{proof}

\begin{remark}[Full-tuple admissibility, Lean-verified]
The full budget-$\le 3$ tuples, of sizes $13,17,15,15,17$ respectively,
are admissible at every prime~$p$. Small primes ($p$ up to the tuple
length) are discharged by \texttt{native\_decide} per residue; larger
primes follow by the same pigeonhole bound used in the core-sub-tuple
proof (\texttt{numForbiddenAtFull\_le\_length} + a per-residue length
equality). See \texttt{hard\_root\_r\_\{0,1716,4862,16302,17160\}\_full\_admissible\_all}
in \texttt{Erdos647HardRootAdmissibility.lean}. The underlying
coefficient data is in \texttt{data/hard\_root\_forms\_r*.json} and
\texttt{data/admissibility\_new\_hard\_residues.json}.
\end{remark}

%% ====================================================================
\section{Barrier theorem}\label{sec:barrier}
%% ====================================================================

\begin{theorem}[Barrier theorem, Lean-verified including CRT]
For any budgeted affine state with $m$ active clauses and any
finite set $S$ of regular primes each greater than~$m$, there
exists a residue avoiding all clause divisibilities simultaneously.
\end{theorem}

\begin{proof}
Induction on $|S|$ via Finset.induction. At each step,
\texttt{one\_prime\_escape} gives a good residue at the new prime;
CRT-composition (proved via ZMod arithmetic) combines with the
inductive hypothesis.
\end{proof}

%% ====================================================================
\section{Conditional prime-core resolution}\label{sec:conditional}
%% ====================================================================

\begin{theorem}[Conditional, under Dickson]
Under Dickson's conjecture for each of the $5$ specific admissible
tuples, the budget-$\le 3$ prime core is infinitely often solvable.
\end{theorem}

The passage from ``prime core solvable'' to ``full \#647 witness''
requires controlling the infinite family of budget-$\ge 4$ constraints.
Section~\ref{sec:killer_sets} records the empirical content at the
finite end of this family: for every strong-threshold record observed
to date (at any residue), some $k\in\{14,15,16\}\cup\{7,11,13\}$
witnesses an Erd\H{o}s-inequality failure, and the minimum killer
set $S_0=\{14,15,16\}$ covers all $228$ strong-threshold records at
$r=0$. However, Remark~\ref{rem:emp_only} shows that the three
finite-incompatibility programs we tested fail to produce a proof of
this phenomenon; the natural next step is a density-level control on
the budget-$\ge 4$ Dickson-set tail, with an analytic bound on the
tuple-restricted Titchmarsh-divisor problem on the Dickson set as the
nearest proxy target. A proxy bound of this form would be a major
step toward closing the Dickson-to-\#647 gap, but a separate bridge
from the divisor proxy to the density-level statement remains open
(direct Markov/Chebyshev on fixed-threshold divisor moments does not
suffice). We do not establish either bound here and flag both as
principal open questions for future work.

We also record, for context, a separate Dirichlet-$L$
pair-correlation / Stage-3B route pursued in parallel project
notes. That route is not an input to the Dickson-based conditional
statement of this section. In its current form it gives a
conditional asymptotic theorem under GRH for primitive $L(s,\chi)$
with $q\mid \alpha$, together with a $\chi$-twisted
Hardy--Littlewood hypothesis $H_\chi$, and a finite-$T^*$
numerical corollary under an additional bookkeeping gate $B_\chi$.
The strict GRH-only finite-$T^*$ individual-$\chi$ version is
stronger, and the completed recon treats it as effectively red in
the present framework after two explicit failed subroutes.

A positive Dickson-conditional resolution of the hard-tuple prime
core would stand in tension with the \texttt{variants.lim}
intuition of Erd\H{o}s (Section~\ref{sec:problem}).

\begin{proposition}[Dickson--$\#647$ dichotomy, localized]\label{prop:dichotomy}
Assume Dickson's conjecture holds for at least one admissible
hard-tuple tuple at $r\in\mathcal{R}^\star:=\{0,1716,4862,16302,17160\}$.
Then for some such $r$ there exist infinitely many $t\in\mathbb{N}$
with every budget-$\le 3$ prime form at $r$ simultaneously prime;
call such $t$ a \emph{Dickson solution at $r$}. Exactly one of the
following holds along any such sequence of Dickson solutions:
\begin{enumerate}
\item[\emph{(A)}] All but finitely many Dickson solutions $t$ at $r$
  fail some budget-$\ge 4$ constraint at $n=2520(46189t+r)$.
  Consequently $\#647$ has at most finitely many solutions.
\item[\emph{(B)}] Infinitely many Dickson solutions $t$ at $r$
  satisfy all budget constraints simultaneously at
  $n=2520(46189t+r)$. Equivalently, $\#647$ admits infinitely many
  solutions.
\end{enumerate}
\end{proposition}

The dichotomy is a statement about the cardinality of $\#647$'s
solution set. Two structural consequences deserve separate
emphasis.
\emph{(i) Relation to \texttt{variants.lim}.} The Erd\H{o}s
\texttt{variants.lim} strengthening asserts
$\max_{m<n}(\tau(m)+m-n)\to\infty$. Case (B) directly contradicts
\texttt{variants.lim} (infinitely many $\#647$ solutions witness
a bounded supremum at arbitrarily large $n$). Case (A) is
\emph{consistent with but strictly weaker than}
\texttt{variants.lim}: finitely many solutions imply the bound is
achieved only finitely often, but \texttt{variants.lim} further
requires unbounded growth, which (A) alone does not guarantee.
\emph{(ii) Localization.} The reduction does not settle which of
(A), (B) holds; but it pins the unresolved content to a precise
arithmetic phenomenon — the budget-$\ge 4$ behaviour along Dickson
solutions at the $5$ specific hard-tuple tuples — that a
sieve-theoretic or analytic result could in principle address. In
particular, the reduction is not merely a partial resolution: it
localizes the \emph{entire} incompatibility between Dickson-type
conjectures and Erd\H{o}s's \texttt{variants.lim} intuition to
these $5$ specific arithmetic witnesses.
\emph{(iii) Target hypothesis.} A natural paper-primary analytic
target is a density-level control on the budget-$\ge 4$ tail along the
Dickson set. The nearest literature-shaped proxy target is a
\emph{tuple-uniform Titchmarsh-type} control on the Dickson set: an
upper bound for divisor moments on a tail form $L_{r,k_0}$ averaged
over the simultaneous-prime-tuple set
$D_r(X)=\{t: L_{r,k}\,\text{prime for all core }k\}$. Existing project
notes freeze that proxy as a uniform upper bound of the shape
\[
  \sum_{t\le X,\ t\in D_r(X)} \tau(L_{r,k_0}(t))
  \;\ll\; |D_r(X)|\log X
\]
for $k_0$ in the finite tail family. Remark~\ref{rem:track2_absolute}
records the current unconditional state of the art on this proxy and
isolates the two remaining conjectural inputs. The proxy alone would
still not settle the full Track~2 tail statement, because fixed-
threshold Markov/Chebyshev on moments of $\tau(L_{r,k_0}(t))$ does
not directly upper-bound the good-set density on the Dickson set.
At the reference residue $r=0$, the bridge now closes conditionally as recorded
in Remark~\ref{rem:track2_absolute}; for the other four hard residues and for
residue-uniform versions, the direct tail-density statement and this tuple-restricted divisor proxy remain open analytic targets. A natural conceptual framework for the proxy bound would be a prime-core-relative analog of Matthiesen's correlation
estimates for multiplicative functions on affine-linear systems
(Matthiesen, \emph{Correlations of the divisor function},
Proc.\ London Math.\ Soc.\ (3) 104 (2012), 827--858;
\emph{Linear correlations of multiplicative functions},
Proc.\ London Math.\ Soc.\ (3) 121 (2020), 372--425), controlling
divisor moments on the $13$-form Dickson set rather than on the
ambient integers.

\begin{remark}[Absolute-form tuple-restricted $\tau$ bound at $r=0$ (unconditional); conditional outlook]\label{rem:track2_absolute}
At the reference hardest residue $r=0$ (tuple rank~$13$), the
proxy bound of \emph{(iii)} has a clean \emph{unconditional}
form; the bridge from the proxy to a Dickson-set density-level
statement sits inside the same Dickson/Schinzel conjectural frame
used for the dichotomy in Proposition~\ref{prop:dichotomy} and is
flagged here as a \emph{conditional outlook}, not as a result of
this paper. Analogous statements at the other four hard residues
$r\in\{1716,4862,16302,17160\}$ (ranks~$15$ or $17$) are expected
to hold by the same composite but require residue-specific
bookkeeping not executed here; we state this Remark only at $r=0$.

\emph{(a) Unconditional absolute form at $r=0$.} For every tail
offset $k_0$ in the finite budget-$\ge 4$ family,
\[
  \sum_{t\le X,\,t\in D_0(X)} \tau(L_{0,k_0}(t))
  \;\ll_{k_0}\; X\,(\log X)^{-12}
\]
unconditionally. The proof is a composition of known published
tools at rank~$13$: a Selberg $\Lambda^{2}$ upper-bound sieve on
the $13$-form Dickson indicator $\mathbf{1}_{D_0(X)}$ (in the
Maynard 2015 normalization; the $\Lambda^{2}$ weight majorizes
$\mathbf{1}_{D_0(X)}$ pointwise on $W$-tricked residues at
admissibility, which is verified for $r=0$ at $p\le 19$ via the
factorization $46189=11\cdot 13\cdot 17\cdot 19$ and the
divisibility $7\mid d_k$ for every $d_k$ in the $r=0$ divisor
set); CRT reduction of the joint divisibility constraints to a
single arithmetic progression via the $W$-trick (Maynard 2015
equations (5.1)--(5.2)), with the cross-form resultant prime
support $\{2,3,5,7,11,13,17,19,23\}$ absorbed at $w(N)\ge 23$;
and Shiu/Nair--Tenenbaum Brun--Titchmarsh for $\tau$ on a linear
form in an arithmetic progression (Shiu, \emph{A Brun--Titchmarsh
theorem for multiplicative functions}, J.\ reine angew.\ Math.\
313 (1980), 161--170; Nair and Tenenbaum, \emph{Short sums of
certain arithmetic functions}, Acta Math.\ 180 (1998), 119--144;
Henriot, \emph{Nair--Tenenbaum bounds uniform with respect to the
discriminant}, Math.\ Proc.\ Camb.\ Phil.\ Soc.\ 152 (2012),
405--424). The rank-$13$ sieve integral is tabulated in the
Polymath8b range
(\emph{Variants of the Selberg sieve, and bounded intervals
containing many primes}, Res.\ Math.\ Sci.\ 1 (2014)).

\emph{(b) Conditional outlook: Dickson-set density at $r=0$.}
Converting the unconditional absolute bound in~(a) into a statement
about the density of the $\tau$-exceptional set \emph{inside} the
rank-$13$ Dickson set $D_0(X)$ requires a quantitative lower bound
of the shape $|D_0(X)|\gg X(\log X)^{-13}$. This is the quantitative
form of the rank-$13$ case of Dickson's conjecture (equivalently, a
specialisation of Schinzel's Hypothesis~H at the $r=0$ $13$-tuple),
and it sits within the same named-conjecture framework already
invoked by Erd\H{o}s himself as the conditional premise under which
the infinitude-variant adjacent to \#647 would follow~\cite{Bloom647}.
Proposition~\ref{prop:dichotomy} already conditions the Track-2
dichotomy on a qualitative form of Dickson (existence of infinitely
many $t\in D_0(X)$); the quantitative lower bound above is a
strictly stronger hypothesis in the same conjectural frame.

Under this quantitative Dickson-type hypothesis, a Dickson-set
Tur\'an--Kubilius concentration bound for $\omega(L_{0,k_0}(t))$
on $D_0(X)$ --- specifically, a second-moment Selberg variance
estimate
$\sum_{t\in D_0(X)} (\omega(L_{0,k_0}(t))-\log\log X)^{2}
 \ll |D_0(X)|\,\log\log X$
--- implies by Chebyshev that
$\{t\in D_0(X):\omega(L_{0,k_0}(t))\le C\}$ has Dickson-set
density $o(1)$ for every fixed $C$, hence the $r=0$ Track-$2$
conclusion $E_0(X)/|D_0(X)|=o(1)$. This is a conditional corollary
of the reduction, not a theorem of the present paper; its status
and the conditional framework are recorded here to make the
analytic content explicit and to distinguish it from the
unconditional proxy in~(a).
\end{remark}

%% ====================================================================
\section{Quantitative predictions and witness search}\label{sec:BH}
%% ====================================================================

The Bateman--Horn singular series for the $13$-tuple at $r=0$ is
$C_{\mathrm{BH}}\approx 2.3179\times 10^8$ (primes $\le 500{,}000$;
convergence ratio $0.999949$).

Using the leading-order Bateman--Horn asymptotic, the expected
witness count is
\[
\#\{t\le X:\text{all 13 prime}\}
  \;\approx\; \frac{C_{\mathrm{BH}}\cdot X}
  {\prod_{i=1}^{13}\log(\alpha_i X)}\,.
\]
At $X=10^{13}$: ${\approx}\,0.48$ expected witnesses.
At $X=10^{14}$: ${\approx}\,2.5$ expected witnesses. The figures
reported are the leading-order Bateman--Horn prediction assuming full
independence across the 13 forms; refinements that account for
residue-dependent correlation among the forms (specifically, the
common divisibility structure induced by $d_i \mid 2520$) typically
reduce the leading-order prediction by $30$--$40\%$ in comparable
settings, but we do not attempt a full correlation-corrected prediction
here. Recomputable from \texttt{scripts/paper\_numerical\_audit.py}.

A computational search is in progress at $r=0$. Phase~1
($t\le 10^{13}$) and two parallel Phase~2 ranges partitioning
$[10^{13}, 10^{14}]$ run concurrently; interim results are
synced to the public repository and refreshed regularly.

\paragraph{Near-miss distribution (snapshot 2026-04-17).}
At this snapshot (Phase~1 at $\sim 17\%$; Phase~2 ranges at
$\sim 2\%$ each), the combined search has observed $9$ distinct
$12/13$ near-misses, $129$ at $11/13$, and $1{,}355$ at $10/13$;
no $13/13$ witness has been found. The $9$ near-misses at $12/13$
span $t = 8{,}618{,}618{,}190$ (the smallest, with composite form
$d=210$, composite factor $1109$, above sieve limit) up to
$t = 55{,}386{,}505{,}208{,}280$.

Across the full corpus of $40{,}984$ composite form-values observed
in the $10{,}526$ near-miss records (at $k\in\{9,10,11,12\}$), the
per-form composite rate is essentially uniform: $19.6\%$ to $20.8\%$
across the $13$ forms, with no form significantly over- or
under-represented. At the extreme $12/13$ tail ($N=9$), the
composite is distributed as $2$ hits each at $d\in\{1260, 2520\}$
(the two largest) and $0$ hits at $d\in\{105, 140\}$ (the two
smallest); we flag this as a preliminary observation pending
larger $N$ from continued search, noting that a $4/9$ concentration
in the two largest of thirteen forms is not statistically distinguishable
from the uniform null at this sample size. The broader uniformity
across $40{,}984$ form-values is the primary observation; the $N=9$
tail pattern is reported as empirical data but not as a structural
claim.

Under naive Bateman--Horn the expected number of $13/13$
witnesses through $X=10^{14}$ is $\approx 2.5$. The current
corpus reaches $X\approx 1.01\times 10^{14}$ and contains exactly
one $13/13$ hit (at $t=11{,}741{,}230{,}040{,}310$); see
Section~\ref{sec:killer_sets} and
Remark~\ref{rem:emp_only} for the empirical context. The
observed-to-predicted ratio at that horizon is $\approx 0.62$,
consistent with BH leading-order density to within a factor of
two. Results are recomputable from
\texttt{data/all\_phases\_merged.jsonl} and
\texttt{docs/near\_miss\_snapshot\_20260417.md}; the full
quantitative refresh is in
\texttt{docs/d2\_bh\_validation\_result.md}.

\begin{proposition}[Per-residue Bateman--Horn predictions]\label{prop:perBH}
The leading-order Bateman--Horn predictions across the $5$ hard
residues differ substantially due to tuple-size variation:

\begin{center}
\begin{tabular}{cccccc}
\hline
$r$ & $K$ & $C_{\mathrm{BH}}$ &
  pred.\ at $10^{13}$ & pred.\ at $10^{14}$ & pred.\ at $10^{16}$ \\
\hline
$0$      & $13$ & $2.32\times 10^{8}$  & $0.47$              & $2.54$               & $78.68$ \\
$1716$   & $17$ & $6.83\times 10^{10}$ & $2\times 10^{-5}$   & $1.4\times 10^{-4}$  & $3.0\times 10^{-3}$ \\
$4862$   & $15$ & $2.53\times 10^{9}$  & $2.6\times 10^{-3}$ & $1.2\times 10^{-2}$  & $3.2\times 10^{-1}$ \\
$16302$  & $15$ & $5.42\times 10^{9}$  & $5.1\times 10^{-3}$ & $2.5\times 10^{-2}$  & $6.4\times 10^{-1}$ \\
$17160$  & $17$ & $1.15\times 10^{11}$ & $4\times 10^{-5}$   & $2.2\times 10^{-4}$  & $5.1\times 10^{-3}$ \\
\hline
\end{tabular}
\end{center}

Computed at prime bound $10^{6}$ via
\texttt{scripts/compute\_bh\_all\_residues.py}; full data in
\texttt{data/bh\_all\_residues.json}.
\end{proposition}

\begin{proposition}[Case-A screen density: pointwise pass probability and cumulative expected count]\label{prop:case_a_density}
In this proposition, \emph{case-A} refers to the fixed-offset
truncated screen $\tau(2520N - k) \le k+2$ at every
$k\in[1,100]\setminus K_{\mathrm{base-enriched}}(r)$.  Any true
Erd\H{o}s~$\#647$ counterexample on residue $r$ must satisfy this
screen as a \emph{necessary} condition (the full problem demands
$\tau(n-k)\le k+2$ for all $1\le k<n$, a condition with $k$ growing
with $n$).  We do \emph{not} claim that the truncated screen alone
implies \#647-finiteness; it is the finite-horizon instrument used to
bound what modeling can currently verify.  The offsets $k>100$ and
the moving-tail regime $k\to\infty$ are not treated in this
proposition.  A moving-$K$ window scan (see
Remark~\ref{rem:moving_k_empirical} and
\texttt{scripts/moving\_k\_window\_experiment.py}) shows that for
$k\in[101, 1000]$, per-$k$ empirical pass rates are $\ge 0.50$ at
every residue with means converging to $1$ as $K$ grows (mean
$0.96$ at $[101, 200]$, $0.99$ at $[201, 500]$, $0.999$ at
$[501, 1000]$).  No very-tight offsets ($<0.3$) appear at $k>100$;
the tight pressure is concentrated at $k\le 100$.

We separate two distinct quantities that were conflated in an earlier
draft.

\emph{(A) Pointwise omitted-offset pass probability $p_r(t)$.}  The
probability that $\tau(2520N - k)\le k+2$ holds at every
$k\in[1,100]\setminus K_{\mathrm{base-enriched}}(r)$ for $t$ sampled
uniformly at scale $t$.  Directly measured at $t\sim 10^{8.5}$ via
per-$k$-independence product (calibration (ii) below), extrapolated to
other scales via per-$B^{\mathrm{eff}}$ pass-rate drift observed in
the scale-stability experiment.  This quantity \emph{decreases} with $t$
(primes and low-$\tau$ integers become rarer as affine forms grow).

\emph{(B) Expected cumulative case-A count $E_r([T_0, X])$.}
\[
  E_r([T_0, X]) \;\approx\; \int_{T_0}^X \rho_r(t)\, p_r(t)\,dt,
\]
where $\rho_r(t) = C_{\mathrm{BH}}/(\log(\alpha\cdot t))^{m_r}$ is the
Bateman--Horn base-tuple hit-density at scale $t$ and $T_0$ is the
lower integration limit.  We report results with $T_0 = 10^8$ (the
floor of the scale-stability calibration range); integrating from
$T_0 = 2$ gives values $10^2$--$10^6$ times larger per residue but
relies on the scale model extrapolated below its fitted range, so we
do not treat it as load-bearing.  The horizon $X$ is the $t$-horizon
throughout this proposition ($X = X_t$), not the $N$- or $n$-horizon.
This quantity is non-decreasing in $X$.

\emph{Two active calibrations plus a legacy heuristic.} Each is
computational or heuristic; none is a proved density theorem.
Calibrations (i) and (ii) are active; (iii) is a retired
historical placeholder retained only for transparency.
\begin{enumerate}
\item[(i.a)] \emph{Raw random-$t$ empirical upper bound.}  A case-A
pass scan at $t\in[10^8, 10^9]$ with $n=1{,}000$ samples per
residue (5{,}000 pooled) gives $0/1000$ per residue, $0/5000$
pooled.  One-sided $95\%$ Clopper--Pearson on the raw IID sample:
$P(\text{full pass})\le 1-0.05^{1/1000}\approx 3.0\times 10^{-3}$
per residue, pooled $\le 1-0.05^{1/5000}\approx 6.0\times 10^{-4}$.
This is a statistical confidence bound on the raw-$t$ distribution
at $t\sim 10^{8.5}$.

\item[(i.b)] \emph{Recorded near-miss corpus diagnostic.}  The
full $r=0$ witness-search corpus recorded to $t\sim 10^{14}$
contains $26{,}253$ logged near-miss records (all primes $\ge 9$
enriched forms).  All $26{,}253$ fail case-A; in particular, the
single $13/13$ shape-pass (the confirmed witness at
$t=11{,}741{,}230{,}040{,}310$) fails case-A at $k=7$ because
$\tau(2520N - 7) = 16 > 9 = k+2$ (direct computation).
Interpreting the corpus as a sample from the recorded-near-miss
distribution, one-sided $95\%$ Clopper--Pearson upper bounds per
stratum are:

\begin{center}
\begin{tabular}{c|cc}
\hline
stratum & sample $n$ & UB on $P(\text{full pass}\mid\text{stratum})$ \\
\hline
$C_{\ge 9}$  & $26{,}253$ & $1.1\times 10^{-4}$ \\
$C_{\ge 11}$ & $330$      & $9.0\times 10^{-3}$ \\
$C_{\ge 12}$ & $21$       & $0.133$ \\
$C_{\ge 13}$ & $1$        & $0.951$ \\
\hline
\end{tabular}
\end{center}

\noindent The $C_{\ge 13}$ row is a single observation and is
diagnostic rather than estimatory: the sole $13$-tuple shape-pass
fails case-A at $k=7$, falsifying the implicit assumption that
shape-pass implies near-full-pass, without yielding a useful
conditional probability estimate.  The corpus-wide $C_{\ge 9}$
bound is a statistical statement about the \emph{recorded-near-miss
layer}, not the raw-$t$ distribution; combining with (i.a) for a
full per-$t$ bound requires an explicit model of the recorded
layer's density within the search range.

\item[(ii)] \emph{Atlas-derived independence-product estimate
(pointwise).}  From $n=2{,}000$-sample per-$k$ pass rates at
$t\in[10^8,10^9]$, the per-$k$-independence product evaluates to:
$3.11\times 10^{-10}$ ($r=0$), $1.66\times 10^{-10}$ ($r=1716$),
$1.16\times 10^{-10}$ ($r=4862$), $1.92\times 10^{-10}$ ($r=16302$),
$1.60\times 10^{-10}$ ($r=17160$).  These are values of $p_r$ at
$t\sim 10^{8.5}$.  The scale-stability experiment at three $t$-scales
spanning four orders of magnitude fits
$p_r(t) \approx p_r(10^{8.5})\cdot 10^{\gamma_r\cdot(\log_{10} t - 8.5)}$
with per-residue slopes $\gamma_r\in[-3.4,-2.0]$.

\item[(iii)] \emph{Legacy witness-scale heuristic (not a measurement).}
Earlier drafts used $P(\text{full case-A pass}\mid \text{shape hit})
\approx 10^{-14}$ as a calibration anchor, rationalized by dividing
the one confirmed shape hit by the search range $\sim 10^{14}$.  The
per-$t$ verification in (i.b) shows that shape-pass does not imply
near-full-pass: the witness fails at $k=7$.  The $10^{-14}$ figure is
therefore a legacy order-of-magnitude placeholder, not a direct
measurement; we retain it only as historical context and do not use
it in the quantitative synthesis below.
Combined with $\rho_r(t)$, yields a cumulative estimate
$E_r(X)\approx 10^{-14}\cdot N_{\mathrm{BH}, r}(X)$ where
$N_{\mathrm{BH},r}$ is the expected shape-hit count.
\end{enumerate}

\emph{Table of pointwise probabilities $p_r(t)$} (calibration (ii) with
scale extrapolation):

\begin{center}
\begin{tabular}{c|ccccc}
\hline
$\log_{10} t$ & $r=0$ & $r=1716$ & $r=4862$ & $r=16302$ & $r=17160$ \\
\hline
$8.5$  & $3.1\times 10^{-10}$ & $1.7\times 10^{-10}$ & $1.2\times 10^{-10}$ & $1.9\times 10^{-10}$ & $1.6\times 10^{-10}$ \\
$10.5$ & $5.4\times 10^{-12}$ & $2.0\times 10^{-13}$ & $1.7\times 10^{-13}$ & $8.8\times 10^{-13}$ & $1.8\times 10^{-13}$ \\
$12.5$ & $9.3\times 10^{-14}$ & $2.5\times 10^{-16}$ & $2.5\times 10^{-16}$ & $4.1\times 10^{-15}$ & $2.1\times 10^{-16}$ \\
$14.5$ & $1.6\times 10^{-15}$ & $3.1\times 10^{-19}$ & $3.8\times 10^{-19}$ & $1.9\times 10^{-17}$ & $2.3\times 10^{-19}$ \\
\hline
\end{tabular}
\end{center}

\noindent\emph{Table of cumulative expected counts $E_r([10^8, X])$}
(integral density model, calibration (ii); integration floor
$T_0 = 10^8$ at the scale-stability calibration lower bound):

\begin{center}
\begin{tabular}{c|ccccc}
\hline
$X$ & $r=0$ & $r=1716$ & $r=4862$ & $r=16302$ & $r=17160$ \\
\hline
$10^{14}$ & $1.2\times 10^{-13}$ & $2.3\times 10^{-17}$ & $6.0\times 10^{-16}$ & $1.1\times 10^{-15}$ & $2.2\times 10^{-17}$ \\
$10^{17}$ & $1.3\times 10^{-13}$ & $2.3\times 10^{-17}$ & $6.0\times 10^{-16}$ & $1.1\times 10^{-15}$ & $2.2\times 10^{-17}$ \\
$10^{19}$ & $1.3\times 10^{-13}$ & $2.3\times 10^{-17}$ & $6.0\times 10^{-16}$ & $1.1\times 10^{-15}$ & $2.2\times 10^{-17}$ \\
\hline
\end{tabular}
\end{center}

\noindent Most residues saturate quickly within the calibration range
because $p_r(t)$ decays faster than $\rho_r(t)$ grows (per-residue
slope $\gamma_r\in[-3.4,-2.0]$). At $r=0$ the slope $\gamma_0 = -2.03$
is shallowest and $E_0([10^8,X])$ grows modestly from
$1.0\times 10^{-13}$ at $X=10^{12}$ to $1.3\times 10^{-13}$ at
$X=10^{19}$.  This saturation is the sharpest quantitative
consequence of the density model: for $r\in\{1716, 4862, 16302,
17160\}$ the integrated expected counterexample count
$E_r([10^8, X])$ is essentially constant for all $X\ge 10^{13}$,
and for $r=0$ it grows by at most $\sim 30\%$ between $X=10^{12}$
and $X=10^{19}$.  Under the current calibration the integrated
case-(A) counterexample risk is therefore approximately
scale-independent at large $X_t$, with total expected count bounded
above by $1.3\times 10^{-13}$ (or $8.1\times 10^{-13}$ under the
anchored log-law) at $r=0$ and by smaller values at the other four
residues, across all $X_t$ in the modeled range.

\emph{Sensitivity envelope} (integral under $\gamma_r\pm 0.5$) at
$X=10^{19}$: for $r=0$ with $\gamma_0=-2.03$,
$E_0([10^8, 10^{19}])$ evaluates to $1.0\times 10^{-12}$ at
$\gamma+0.5$ (less-negative slope, slower decay, larger tail
integral), $1.3\times 10^{-13}$ at the fitted value, and
$7.6\times 10^{-14}$ at $\gamma-0.5$ (more-negative slope, faster
decay, smaller integral).  Total range factor $\sim 14$; other
residues vary by $<3\times$ under the same envelope.

\emph{Calibration-floor caveat.} Integrating from $T_0 = 2$ rather
than $T_0 = 10^8$ yields values $10^2$--$10^6$ times larger per
residue, but relies on the scale model extrapolated below its fitted
range and is not load-bearing.  The $[T_0, X]$ framing is the
credible form.

\emph{Comparison with log-law tail model and witness-calibrated (iii).}
Calibration family (ii) has two variants: (ii.a) the local power-law
fit $p_r(t)\propto t^{\gamma_r}$, and (ii.b) an anchored Landau-Sathe
log-law tail $F_B(x)\propto (\log\log x)^{s-1}/(\log x)$.  The
Landau model is calibrated from the scale-stability data for the
decay \emph{shape} and anchor-normalized at $t_0 = 10^{8.5}$ to the
empirical independence product (correcting a factor-$\sim 6$--$653$
per-residue normalization offset in the raw Landau fit).  At
$X=10^{19}$ the cumulative estimates are:

\begin{center}
\begin{tabular}{c|ccc}
\hline
$r$ & (ii.a) power-law & (ii.b) anchored log-law & (iii) witness-calibrated \\
\hline
$0$     & $1.3\times 10^{-13}$ & $8.1\times 10^{-13}$ & $3.3\times 10^{-9}$ \\
$1716$  & $2.3\times 10^{-17}$ & $2.7\times 10^{-16}$ & $1.7\times 10^{-12}$ \\
$4862$  & $6.0\times 10^{-16}$ & $7.7\times 10^{-15}$ & $7.6\times 10^{-11}$ \\
$16302$ & $1.1\times 10^{-15}$ & $1.1\times 10^{-14}$ & $7.6\times 10^{-11}$ \\
$17160$ & $2.2\times 10^{-17}$ & $2.8\times 10^{-16}$ & $1.7\times 10^{-12}$ \\
\hline
\end{tabular}
\end{center}

\noindent The anchored log-law (ii.b) gives values $6$--$13\times$
larger than the power-law (ii.a) per residue, consistent with the
Landau asymptotic decaying more slowly than a local power-law fit.
Both variants of family (ii) predict negligible expected counts at
the modeled horizons, and are consistent with the empirical upper
bounds of (i.a) and (i.b).  The legacy heuristic (iii) gives larger
values ($10^3$--$10^4\times$ compared with (ii)), but the witness
verification in (i.b) shows (iii) was never a direct measurement;
its numerical value is retained only for historical comparison.
See \texttt{scripts/loglaw\_tail\_model.py} and
\texttt{data/loglaw\_tail\_model.json}.

\emph{Modeling caveats.}  Four specific sources contribute to the
gap between (ii) and (iii): (a) the independence-product assumption
in (ii) (cross-$k$ correlations may change (ii) by several orders of
magnitude in either direction; for omitted offsets $n-a,n-b$ with
primes $q>|a-b|$ the two divisibility indicators cannot simultaneously
vanish, giving a negative covariance at each such prime, while local
Bateman--Horn factors can contribute positively — the sign and size of
the net effect on joint low-$\tau$ pass events is not yet modeled);
(b) the marginal-over-$t$ versus per-shape-hit framework mismatch;
(c) the single-data-point variance in (iii); and (d) the power-law
scale extrapolation $p_r(t)\propto t^{\gamma_r}$ used in (ii).
Regarding (d): per-$B^{\mathrm{eff}}$ pass rates are asymptotically
governed by Landau / PNT-type log-law probabilities
$F_b(Z)\propto (\log\log Z)^{b-2}/(\log Z)$, not by power laws in $t$.
The anchored log-law variant (ii.b) above uses the Landau-Sathe decay
shape anchored to the empirical independence product at $t=10^{8.5}$.
At $X=10^{19}$ the anchored log-law gives values $6$--$13\times$
larger than the power-law, consistent with a slower-decaying
asymptotic.  This is a modest correction relative to the overall
$10^3$--$10^4\times$ gap between family (ii) and the witness
heuristic (iii): the tail-law choice is not a first-order source of
the remaining uncertainty.  We consider caveat (d) substantially
tested and demoted.

Regarding (a): we tested the independence-product assumption by
computing direct block ratios
$R_S = \Pr(\bigcap_{k\in S}\mathrm{Pass}_k) / \prod_{k\in S}\Pr(\mathrm{Pass}_k)$
for four block families at each hard residue: the monitor
$\{7,11,13,16\}$, the empirical minimum cover $\{15,16,27\}$, the top-$5$
bottleneck offsets from calibration (ii.a), and the union of monitor
and top-$5$ bottlenecks.  At $n = 5{,}000$ samples per residue and
$t\in[10^8,10^9]$, all twenty $R_S$ values fall in $[0.63, 1.80]$, with
$\max|\log R_S|\approx 0.59$.  Independence holds empirically to
within factor $\sim 2$ at the block level.  The $10^3$--$10^4\times$
gap between (ii) and (iii) is therefore \emph{not} attributable to
cross-$k$ dependence; caveat (a) is substantially tested and
demoted.

The remaining primary source of modeling uncertainty is caveat (b),
the marginal-over-raw-$t$ versus per-shape-hit conditioning mismatch.
Calibration (ii) samples $t$ uniformly while calibration (iii) is
conditional on a base-tuple prime hit.  A weighted base-tuple
conditional surrogate is the next natural refinement.
Neither calibration is a proved bound; all are extrapolated models
with documented uncertainties.

See \path{scripts/density_integral_model.py},
\path{data/density_integral_model.json},
\path{scripts/scale_stability_test.py}, and
\path{data/scale_stability_test.json}.

The largest single-$k$ bottleneck in (ii) is $k=16$ (per-$k$ pass
rate $\approx 0.35$ at every residue), consistent with
Proposition~\ref{prop:k16_universal}.  See
\path{scripts/case_a_density_factored.py} and
\path{data/case_a_density_factored.json} for the full per-$k$
rate tables; \path{docs/case_a_density_model.md} for caveats.
\end{proposition}

\begin{remark}[Tuple-size convention and divisor-extension sizes]\label{rem:tupleconv}
Proposition~\ref{prop:perBH} uses the hard-root tuple sizes
$K\in\{13,17,15,15,17\}$ at which the singular series were
originally computed. The divisor-extension tuple sizes from
Section~\ref{sec:div_ext} and Theorem~\ref{thm:complete} are
$\{13,15,14,14,15\}$ at $r\in\{0,1716,4862,16302,17160\}$. The two
conventions agree at $r=0$, so the r=0 entry of the table is
directly comparable to the divisor-extension framework used
elsewhere in the paper. At nonzero residues, the divisor-extension
convention has smaller $K$ by $1$ or $2$; at $X=10^{14}$ this
shifts the predicted density upward by a factor of
$\log(\alpha X)^{K_\text{old}-K_\text{new}}\approx 10^{1.5}$ to
$10^{3}$, which does not change the order-of-magnitude picture
(nonzero residues remain effectively zero at current search horizons).
Proper recomputation of the singular series on the
divisor-extension tuples is deferred to follow-up work; for the
paper's structural conclusions (the empirical killer-set
discussion of Section~\ref{sec:killer_sets} and the conditional
resolution of Section~\ref{sec:conditional}) the effect is
within existing order-of-magnitude statements.
\end{remark}

\begin{remark}[Informational focus of the $r=0$ search]
The $r=0$ residue is by far the most informative computational
target. At $X=10^{14}$, the two $17$-form residues ($r\in\{1716,
17160\}$) have expected witness counts roughly $10^{4}$ below
$r=0$, while the two $15$-form residues ($r\in\{4862,16302\}$) are
roughly $10^{2}$ below $r=0$. The separation arises from the
tuple-size factor in the denominator product
$\prod\log(\alpha_i X)$: each additional form incurs a $\log X$
penalty, and at $\log 10^{14}\approx 32$, even two extra forms
reduce the predicted density by a factor of $\approx 10^{3}$. The
asymmetry is robust in its direction to refined-formula corrections
(correlation adjustments to the singular series typically move
predictions by factors of a few to at most an order of magnitude,
not by the four orders separating the $17$-form residues from
$r=0$); we have not, however, computed the refined predictions for
the non-$r=0$ residues.

The practical consequence: a null result from the witness search at
$r=0$ (Phases~1 and~2) is informative — at $X=10^{14}$ the
probability of observing zero witnesses under leading-order
Bateman--Horn is $e^{-2.5}\approx 8\%$. Null results at the other
four hard residues at any feasible $X\le 10^{14}$ are, by contrast,
fully consistent with Bateman--Horn and carry little evidence
against witness existence. The search strategy's focus on $r=0$
reflects this informational asymmetry; extending the search to
$r\in\{1716,4862,16302,17160\}$ would not be a productive use of
computational resources in the near term.
\end{remark}

%% ====================================================================
\section{The reduction chain}\label{sec:chain}
%% ====================================================================

The Lean formalization packages the large-$n$ conclusion as a pair of
theorems. The conditional reduction
\[
\texttt{erdos647\_conditional\_on\_openResiduesStage1}:\;
\forall\, n,\; 24<n \;\to\; \lnot\,\texttt{Candidate}\;n,
\]
quantified over any hypothesis asserting that no $N$ with
$N\bmod M\in\texttt{openResiduesStage1}$ yields an Erd\H{o}s~\#647
candidate, captures the entire structural reduction.
The full (unquantified) theorem
\[
\texttt{erdos647\_with\_stage1\_axiom}:\;
\forall\, n,\; 24<n \;\to\; \lnot\,\texttt{Candidate}\;n
\]
applies the single intentional Stage-1 axiom
\texttt{openResiduesStage1\_no\_solution} to the conditional form.
Its current dependency profile is:

\begin{center}
\begin{tabular}{@{}>{\raggedright\arraybackslash}p{0.30\textwidth}
                >{\raggedright\arraybackslash}p{0.45\textwidth}
                >{\raggedright\arraybackslash}p{0.18\textwidth}@{}}
\hline
\textbf{Step} & \textbf{Theorem} & \textbf{Status} \\
\hline
Finite coverage ($25\le n\le 84$)
  & \texttt{candidate\_false\_finite} & proved (\texttt{native\_decide}) \\
Candidate $\to$ IsErdos647
  & \texttt{bridge\_candidate\_to\_isErdos647} & proved \\
$2520$-divisibility
  & \texttt{erdos647\_div\_2520} & proved \\
Sieve reduction (Bridge~2)
  & \texttt{bridge\_isErdos647\_to\_sieve} & proved \\
$55$-residue closure
  & $\texttt{singleOverlapResidues42}\cup
     \texttt{directFullValueResidues13}$ & proved \\
Admissibility ($\times 5$)
  & \texttt{hard\_root\_r\_*\_admissible\_all} & proved \\
Open residues ($\times 41$)
  & \texttt{openResiduesStage1\_no\_solution} & axiom \\
\hline
\end{tabular}
\end{center}

No source-level \texttt{sorry} remains in \texttt{lean/}. The prior
Bridge~3 budget-transfer lane has been retired (the peeled-cofactor
implication was shown to be pointwise unsound; see
\texttt{docs/bridge3\_design\_memo\_20260421.md}) and
\texttt{Bridge3.lean} is now a $16$-line placeholder stub. The
operative closure surface is the Stage-1 partition
\begin{align*}
\texttt{closedResiduesStage1}
&= \texttt{singleOverlapResidues42}\cup\texttt{directFullValueResidues13},\\
\texttt{openResiduesStage1}
&= \texttt{survivingResidues96}\setminus\texttt{closedResiduesStage1},
\end{align*}
with \texttt{closedResiduesStage1\_card = 55} and
\texttt{openResiduesStage1\_card = 41}, both proved by
\texttt{native\_decide}. Bridge~2 is closed in \texttt{Bridge2.lean};
the substantive sub-lemma \texttt{tau\_ge\_large\_noncoprime} is proved
across all six non-trivial coefficient values.

The narrow axiom \texttt{openResiduesStage1\_no\_solution} is confined
to the $41$ residues in the exact-v2 open set and represents the
genuinely open content of Problem~\#647. It subsumes both the
historical hard-$14$ core (five original plus nine exact-v2-reopened
residues) and the twenty-seven residues whose closures under the
invalidated Bridge-3 peeled-cofactor mechanism require a sound
replacement. A post-retirement audit of the $55$ Lean-closed residues
(\texttt{docs/audits/55\_closed\_post\_bridge3\_audit\_20260424.md})
confirms none of the $55$ depends on the Bridge-3 transfer; the
single-overlap and direct full-value arguments work directly at the
full-value level $2520\cdot N-k$ and rely only on divisor
monotonicity and sound coprime-multiplicativity with explicit
coprimality hypotheses.

\begin{remark}[$v_{23}$ Hensel-rootline uniformity of the
sieve-surviving residues]\label{rem:v23_uniformity}
The $55$-closed / $41$-open partition in the Stage-$1$ axiom reflects
the structural reach of this paper's closure mechanisms
(\texttt{SingleOverlapClosure}, \texttt{DirectFullValueClosure}),
not a fundamental hardness gradient between the two classes. A
uniform cycle-$3$-style partial-closure mechanism applies to every
one of the $96$ sieve survivors via the prime~$23$: for each
residue $r\in\texttt{survivingResidues96}$, there exist
shift $k\in\{1,\dots,95\}$ and a sub-AP
$\{N\equiv b_r\pmod{23^3\cdot M}\}$ of density $23^{-3}=1/12167$
on which $2520\cdot N-k = d(r,k)\cdot(\alpha\cdot s + c_0)$ algebraically
with $c_0\ge 2$ and $\tau(d\cdot p)>k+2$ for every prime $p$, so that
no $N$ on the sub-AP is $\texttt{goodOn}$ at~$k$. The structural
basis is $\gcd(23,2520)=\gcd(23,M)=1$, so $23$ is free of both the
divisor budget $2520=2^3\cdot3^2\cdot5\cdot7$ and the sieve modulus
$M=46189=11\cdot13\cdot17\cdot19$; the $23$-adic Hensel tower is
therefore available at every residue class. An empirical scan
verified positive-margin closures at $136$/$136$ tested residues
(the $41$ open, the $55$ closed, and $40$ sampled non-surviving
residues outside \texttt{survivingResidues96}). Consequently the
present paper's $55$/$41$ split is an artifact of \emph{which}
mechanisms are formalized in the $55$ closures and not a structural
claim about the $41$ being harder; a multi-shift primitive
combining several such per-residue sub-AP closures (design in
\texttt{docs/lean\_design/multishift\_primitive\_design\_20260424.md})
would discharge \texttt{openResiduesStage1\_no\_solution} uniformly.
Per-residue certificate data are tabulated in
\texttt{data/all41\_v23\_hensel\_certificates\_20260424.json};
full findings are recorded in
\texttt{docs/audits/all41\_hensel\_v23\_unified\_20260424.md} and
\texttt{docs/audits/v23\_universal\_pattern\_20260424.md}.
\end{remark}

%% ====================================================================
\section{Formal verification summary}\label{sec:formal}
%% ====================================================================

\begin{center}
\begin{tabular}{lcc}
\hline
\textbf{File} & \textbf{Sorry} & \textbf{Lines} \\
\hline
\texttt{Erdos647Main} & 0 (legacy wrapper retired) & 88 \\
\texttt{Erdos647SieveCertificate} & 0 & 64 \\
\texttt{Erdos647HardRootAdmissibility} & 0 & 314 \\
\texttt{Erdos647Closure92\_Legacy} (historical) & 0 & 94 \\
\texttt{Erdos647BarrierTheorem} & 0 & 252 \\
\texttt{Erdos647BridgeV1} & 0 & 904 \\
\texttt{Erdos647B23Closure} & 0 & 316 \\
\texttt{Erdos647ClosureCerts} & 0 & 267 \\
\texttt{Erdos647ClosureData} & 0 & 2837 \\
\texttt{Erdos647BranchForms} & 0 & 1212 \\
\texttt{Erdos647ReductionChain} & 0 & 118 \\
\texttt{Erdos647ResiduePartitionStage1} & 0 & 124 \\
\texttt{Erdos647Stage1Axiom} & 0 + 1 narrow axiom & 27 \\
\texttt{Erdos647ConditionalFiniteWindow} & 0 & 152 \\
\texttt{Bridge2} & 0 & 1044 \\
\texttt{Bridge3} (stub, retired) & 0 & 16 \\
\hline
\textbf{Total (proof-critical subset)} & \textbf{0 sorry + 1 axiom} & \textbf{7829} \\
\hline
\end{tabular}
\end{center}

All finite-computational claims are discharged via Lean's
\texttt{native\_decide}. The proof-critical chain is conditional on the
Stage-1 axiom wrapper:
\texttt{Erdos647ReductionChain.erdos647\_conditional\_on\_openResiduesStage1}
states the theorem relative to the open residue set, and
\texttt{Erdos647Stage1Axiom.erdos647\_with\_stage1\_axiom} applies the
single intentional axiom
\texttt{openResiduesStage1\_no\_solution} to yield the full result.
Bridge~3 is retired as a stub; no source-level \texttt{sorry} remains
in \texttt{lean/}. The complete Lean development spans $49$ files; the
above table records the proof-critical subset. The non-definitional
dependency profile is:
\texttt{propext}, \texttt{Classical.choice}, \texttt{Quot.sound},
\texttt{Lean.ofReduceBool} (from \texttt{native\_decide}),
\texttt{openResiduesStage1\_no\_solution}.

\paragraph{Trust chain and reproducibility.}
All residue closures, shape decompositions, sieve enumerations,
cofactor factorisations, and arithmetic verifications in this paper
are discharged inside the Lean~4 kernel. Finite computational claims
(residue-list cardinalities, divisibility checks, primality over
fixed small ranges, enumeration of surviving residues mod~$M$) use
Lean's \texttt{native\_decide}, which compiles the relevant
decidable proposition to native code and trusts the result via the
\texttt{Lean.ofReduceBool} axiom listed above. No step of the
verified chain depends on externally-computed data ingested from
outside Lean, on unsafe Lean features (\texttt{unsafe}, \texttt{extern}),
or on trusted code outside the Lean source tree.
The residue-list literals \texttt{closedResiduesStage1},
\texttt{singleOverlapResidues42}, \texttt{directFullValueResidues13},
\texttt{openResiduesStage1}, and \texttt{survivingResidues96} are
Lean-internal \texttt{Finset~\ensuremath{\mathbb{N}}} definitions,
and the partition equalities and cardinality claims are proved by
\texttt{native\_decide} on those definitions.
The post-retirement audit at
\texttt{docs/audits/55\_closed\_post\_bridge3\_audit\_20260424.md}
further verifies that the $55$ closures do not depend on
Bridge-3-adjacent reasoning: all multiplicativity uses in
\texttt{Bridge2} and downstream files carry an explicit
\texttt{Nat.Coprime} hypothesis, so the pointwise-unsound
peeled-cofactor identity never appears in the closure proofs.

\begin{thebibliography}{9}
\bibitem{Bloom647} T.\,F. Bloom, \emph{Erd\H{o}s Problem \#647},
\url{https://www.erdosproblems.com/647}.
\bibitem{BatemanHorn} P.\,T. Bateman, R.\,A. Horn, \emph{A heuristic
asymptotic formula concerning the distribution of prime numbers},
\emph{Math.\ Comp.}\ \textbf{16} (1962), 363--381.
\bibitem{Ingham} A.\,E. Ingham, \emph{Mean-value theorems in the theory
of the Riemann zeta-function}, Proc.\ London Math.\ Soc.\ (1927).
\bibitem{TaoParity} T. Tao, \emph{The parity problem in sieve theory},
What's New blog (2007).
\bibitem{DeepMind} DeepMind, FormalConjectures/ErdosProblems/647.lean,
\url{https://github.com/google-deepmind/formal-conjectures} (specifically the theorem \texttt{erdos\_647.variants.lim} in \texttt{FormalConjectures/ErdosProblems/647.lean}).
\bibitem{Dutta} S.~Dutta, comments on Erd\H{o}s Problem \#647,
\url{https://www.erdosproblems.com/647}: initial derivation at
13:40 on 18~Jan 2026 (modular constraints $n=2520N$ with
$d_iN-1$ prime or almost-prime for $d_i\in\{1260,840,630,420,
315,210,105\}$), and refined form at 17:49 on 19~Jan 2026
($315N-1\in\{p,2p\}$, $105N-1\in\{p,p^{2},2p,4p\}$;
$2520N-1$ upgraded to prime after the mod-$4$ observation).
\bibitem{Alexeev} B.~Alexeev, comment on Erd\H{o}s Problem \#647
at 00:50 on 19~Jan 2026,
\url{https://www.erdosproblems.com/647}: flags coprimality gap in
Dutta's $k\in\{8,24\}$ step, upgrades $2520N-1$ to prime via
$2520N-1\equiv 3\pmod 4$, conjectures $11\mid n$ and $13\mid n$,
and exhibits the explicit near-witness
$n = 604{,}517{,}614{,}941{,}240$ (with $N=239{,}887{,}942{,}437$)
for which $\tau(n-k)\le k+2$ holds for $1\le k\le 13$ and $k=24$.
\bibitem{Agbanwa} J.~Agbanwa, submission on Erd\H{o}s Problem \#647,
\url{https://www.erdosproblems.com/647} forum comment at 15:18 on
28~Jan 2026, with companion Lean formalization at
\url{https://zenodo.org/records/18390414}. Flagged as incorrect in
subsequent forum review.
\bibitem{TaoForum647} T.~Tao, comments on Erd\H{o}s Problem \#647,
\url{https://www.erdosproblems.com/647}: forum comment at 22:41 on
2~Oct 2025 (curated description update on the $\tau/\omega$
hierarchy) and at 16:10 on 28~Jan 2026 (response to~\cite{Agbanwa}).
\bibitem{TaoTeravainen2025} T.~Tao, J.~Ter\"av\"ainen,
\emph{Quantitative correlations and some problems on prime factors
of consecutive integers}, preprint arXiv:2512.01739 (1 Dec 2025).
Establishes a general two-point correlation estimate for bounded
multiplicative functions with logarithmic savings (via Pilatte
correlations), a Maynard-type sieve argument for
$\omega(n+k)\le\Omega(n+k)\ll k$ infinitely often (resolving Erd\H{o}s
\#248), and an asymptotic for $\#\{n\le x:\omega(n)=\omega(n+1)\}$ at
almost all $x$. Remark~1.2 of the preprint explicitly flags
Erd\H{o}s~\#647 as ``beyond the capability of our methods,'' noting
that the small-$k$ barrier-case variant is of ``comparable difficulty
to the prime tuples conjecture.''
\bibitem{Lau2026} C.\,F. Lau, \emph{On the Number of Prime Factors of
Consecutive Integers}, preprint arXiv:2604.15042 (16 Apr 2026).
Refines~\cite{TaoTeravainen2025} to $\omega(n+k)\ll\log k$ for
$k\ge 2$ on an infinite sequence; cited on the curated Bloom~\#647
page~\cite{Bloom647}.
\end{thebibliography}

\end{document}
